\documentclass[a4paper,11pt,openany]{ltjsbook}
\usepackage{luatexja}
\usepackage{luatexja-fontspec}
\usepackage{amsmath,amssymb,amsfonts,amsthm}
\usepackage{geometry}
\geometry{top=24mm,bottom=24mm,left=24mm,right=24mm,footskip=10mm}
\usepackage{tcolorbox}
\tcbuselibrary{skins,breakable}
\usepackage{hyperref}
\usepackage{enumitem}
\usepackage{booktabs}
\usepackage{fancyhdr}

\definecolor{carrollred}{RGB}{120, 25, 35}
\definecolor{carrollgrey}{RGB}{80, 85, 95}
\definecolor{softpink}{RGB}{253, 246, 247}
\definecolor{darkslate}{RGB}{30, 35, 45}
\definecolor{boxframe}{RGB}{140, 50, 60}

\hypersetup{
    colorlinks=true,
    linkcolor=carrollred,
    citecolor=carrollgrey,
    urlcolor=carrollred
}

\pagestyle{fancy}
\fancyhf{}
\fancyhead[LE,RO]{\thepage}
\fancyhead[LO]{\nouppercase{\rightmark}}
\fancyhead[RE]{\nouppercase{\leftmark}}
\renewcommand{\headrulewidth}{0.4pt}

\title{\Huge\textbf{論理の遊戯}\\[4mm]
\huge\textbf{ゲームとダイアグラムで楽しむ三段論法と記号論理}\\[6mm]
\Large\textbf{—— 原著全4章 完全翻訳・数学校閲版 ——}\\[5mm]
\large 原著：ルイス・キャロル（Lewis Carroll / Charles L. Dodgson, 1887）\\[2mm]
\normalsize Project Gutenberg eBook \#4763: \textit{The Game of Logic}}

\author{\textbf{ルイス・キャロル} 著\\[2mm]
\textbf{教育情報分析研究会} 訳・数理校閲}
\date{2026年10月 校訂第1版}

\begin{document}

\maketitle

\frontmatter

\chapter*{解題：ルイス・キャロル『論理の遊戯』の数学的・論理学的意義}
\addcontentsline{toc}{chapter}{解題：ルイス・キャロル『論理の遊戯』の数学的・論理学的意義}

『不思議の国のアリス』の著者として世界中で親しまれているルイス・キャロル（本名：チャールズ・ラトウィッジ・ドジソン, 1832--1898）は、オックスフォード大学クライスト・チャーチで26年間にわたり数学講読講師（Mathematical Lecturer）を務めた本物の数学者・論理学者であった。

彼が1887年に著した『論理の遊戯』（\textit{The Game of Logic}）は、アリストテレス以来の伝統的・形式的論理学（三段論法）が持つ形式的な退屈さと難解さを一掃し、正方形のダイアグラム（盤面）と赤・灰のカウンター（駒）を用いて誰でもパズルのように論理推論を楽しめるように設計された、驚くべき先駆的著作である。

本書でキャロルが考案した二重・三重の正方形ダイアグラムは、ベン図（ジョン・ヴェン）と並び、現代のコンピュータ工学における「カルノー図（Karnaugh Map）」の原型をなすものであり、視覚的に論理の真理値や属性の包含関係を決定する画期的な手法であった。

本書は全4章から構成され、
\begin{enumerate}[label=\textbf{第\arabic*章:}]
    \item \textbf{古きを捨てて新しきを（New Lamps for Old）}：命題の本質、三段論法（Syllogism）、誤謬（Fallacies）の分類、および正方形ダイアグラムを用いた推論の基礎。
    \item \textbf{問いかけの数々（Cross Questions）}：小ダイアグラム・大ダイアグラムを用いた命題の表現と記号解釈の徹底演習問題群。
    \item \textbf{歪んだ答え（Crooked Answers）}：第2章の演習に対する詳細な模範解答と、三段論法の具体的導出プロセス。
    \item \textbf{当たるも八卦当たるぬも八卦（Hit or Miss）}：日々の生活や言語の中に潜む論理的推論と誤謬を見抜く総合論理パズル。
\end{enumerate}
を収載している。

本全訳版は、Project Gutenberg eBook \#4763として公開されている1887年初版テキストを底本とし、ローカル大規模言語モデル（\texttt{qwen2.5-7b-instruct} による忠実な翻訳および \texttt{qwen2.5-math-7b-instruct} による数理・LaTeX校閲）を駆使して完全翻訳・製本したものである。

\tableofcontents

\mainmatter


\chapter{古きを捨てて新しきを（New Lamps for Old）}
\label{ch:i}

\begin{center}
\large\textit{命題論・三段論法・誤謬論と正方形ダイアグラムによる論理表現}
\end{center}
\vspace{4mm}

第一章。

古い灯りに新しい灯りを。

「光来，光去。」

1. 命題。

「新しいケーキは美味しい。」
「新しいケーキは美味しくない。」
「すべての新しいケーキは美味しい。」

これらが三つの「命題」です——このゲームで使用する唯一の三種類です：そして最初に行わなければならないことは、これらの命題をボード上で表現することです。

まず、

「新しいケーキは美味しい。」

について始めましょう。

しかし、その前に、重要な注釈があります——一瞬で全て理解するのは決して簡単ではありませんので、この部分は非常に注意深く読み進めてください。

世界には多くの「事物」（例えば、「パン」、「子供」、「甲虫」、「バット」など）が存在します；これらの事物には多くの「属性」（例えば、「焼かれた」、「美しい」、「黒い」、「壊れた」など：実際には、どんな事物にも「属する」と言える何かはすべて属性です）。私たちが何かの事物を指すときには名詞を使用し、何かの属性を指すときには形容詞を使用します。人々は「何かの事物が属性を持たないで存在できるか？」という質問をしました。これは非常に難解な問いであり、私はそれを完全に解答しようとしないつもりです：私たちはそのような問いを無視して軽蔑的に静止するべきです。しかし、彼らが別の角度から質問し、「何かの属性が所有する事物がないで存在できるか？」と尋ねる場合、すぐに答えることができます。「いいえ：それは子供が誰も見守らないまま鉄道旅行をするよりも不可能でしょう！」あなたは空中に「美しい」という言葉を見たり、床に散らばっている「美しい」という言葉を見たことはありませんか？

長々とした説明の目的は何なのでしょうか。それは、次のようなことです。二つの事物（例えば、「いくつかの豚は太った動物である」）や二つの属性（例えば、「ピンクは薄い赤色である」）の名前間に「は」または「である」という言葉を入れることで、意味が成立します。しかし、事物の名前と属性の名前に「は」または「である」という言葉を入れると（例えば、「いくつかの豚はピンクである」）、意味が成立しない場合があります（なぜなら、事物が属性であるとは言えませんよね？）。ただし、話相手との理解があれば例外となります。そして、最も単純な理解としては、述語が文末に繰り返されることを仮定することでしょう。つまり、完全な文章では「いくつかの豚はピンク（豚）」と書くことになります。このように、「である」という言葉が意味を持つようになります。

したがって、「いくつかの新しいケーキは美味しい」という命題を意味するためには、完全な形で「いくつかの新しいケーキは美味しい（ケーキ）」と仮定しなければなりません。この命題には二つの「項」があります。「新しいケーキ」が一つの項であり、「美味しい（ケーキ）」がもう一つの項です。「新しいケーキ」について話しているので、これは「主語」と呼ばれます。「美味しい（ケーキ）」は「述語」と呼ばれます。また、この命題は「部分的」なものと呼ばれ、全体ではなく一部のみを述べているからです。一方、「全般的」な命題とは、その全体を否定するものや肯定するものを指し、ここでは「新しいケーキの全てが美味しい」という命題と「新しいケーキの全てが美味しいではない」という命題を指します。最後に、「命題」自体の定義について知りたいなら、次のように説明できます。「ある特定のクラス（主語）に属する事物の一部または全部が、別の特定のクラス（述語）に属することを述べる文」です。

これらの七つの言葉—命題、属性、項、主語、述語、部分的、全般的—は、友人が論理について学んだかどうかを尋ねたときに、非常に便利で魅力的なものとなります。すべての七つの言葉をあなたの答えに含めれば、友人は深く感動し、「悲しげで賢い男」と去るでしょう。

現在、ボード上の小さな図を見つめてください。そして、それが世界中のケーキを収納するためのキャビネットであると想像してください（もちろん、それはかなり大きなものでなければなりません）。そして、新しいケーキは上半分（「x」でマーク）に、それ以外のケーキ（つまり古いケーキ）は下半分（「x''」でマーク）に配置しましょう。したがって、下半分には古いケーキ、年季の入ったケーキ、洪水以前のケーキなど（もし存在するなら：私は多く見かけたわけではありません）が含まれるでしょう。また、美味いケーキは左半分（「y」でマーク）に、それ以外のケーキ（つまり味が悪いケーキ）は右半分（「y''」でマーク）に配置しましょう。現在では、「x」を新しいケーキ、「x''」を古いケーキ、「y」を美味しいケーキ、「y''」を美味しいではないケーキと理解する必要があります。

次に、 compartment No. 5にはどのような種類のケーキが存在すると予想されますか？

この compartmentは上半分の一部であり、したがって、中にケーキがある場合でもそれは新しいケーキであることがわかります。また、これは左半分の一部でもありますから、そのケーキは美味しいものであることもわかります。したがって、この compartmentにケーキがある場合、そのケーキには「新しい」および「美味しい」という二重属性「new and nice」を持つことになります：もしくは、文字を使用すると、「x y」となります。

注意していただきたいのは、この compartmentの2つの辺には x, y という文字が書かれていることです。これは、どの compartmentにどのような属性があるかを知るための非常に便利なルールです。例えば、compartment No. 7について考えてみましょう。もしケーキがある場合、それは「x' y」であると予想されます：つまり、「新しいではない」と且つ「美味しい」ということになります。

次に、新たな合意を設けましょう。 compartment内に赤い駒が存在すると、その compartmentは「占領されている」として解釈し、つまりその compartmentには少なくとも一つのケーキがあることを意味します（論理において、「いくつか」は「一つ以上」を指すので、 compartment内の一つのケーキでも十分な理由となります）。また、 compartment内に灰色の駒が存在すると、その compartmentは「空いている」として解釈し、つまりその compartmentにはケーキがないことを意味します。以下の図では、赤い駒を置くべき場所には「1」（「一つ以上」を意味する）、灰色の駒を置くべき場所には「0」（「一つもなし」を意味する）を記すことにします。

当論文の主題は「新しいケーキ」であるため、現在は上段の棚のみに関心を持ちます。そこにはすべてのケーキが属性 x を持つものがあり、「新しい」という意味です。

この上段に注目し、もし次のようなマークを見つけたとしましょう,

-----------
                     |     |     |
                     |  1  |     |
                     |     |     |
                      -----------

つまり、No.5の位置に赤い玉が置かれている場合、これは「新しいケーキ」のクラスについて何を教えてくれるでしょうか？

これにより、「新しい」という属性 x を持つ compartment 内にはいくつかのケーキがあることが分かります。すなわち、「新しい」という属性 x だけでなく、さらに「美味しい」という属性 y も持っているいくつかのケーキが存在することを意味します。これは「一部の x-ケーキは y-(ケーキ)」と表現できますか、または文字で置き換えると

「一部の新しいケーキは美味しい（ケーキ）」

あるいは、より短く表すと

「一部の新しいケーキは美味しい」

と表現できます。

最後に、このセクションの最初の命題をどのように表現するかが分かったのです。もし私が述べたすべてが完全には理解できていない場合は、それを繰り返し読み、理解できるまで進めないでください。その一旦マスターした後では、残りは非常に簡単になります。

他の命題を処理するために少し手間を省くため、一貫して「ケーキ」の言葉を使わないことにしましょう。私は棚が使用されるすべてのもののクラス全体を「宇宙」と呼びます。したがって、私たちはこの業務を始めるために「ケーキの宇宙を取りましょう」というように始めることができます。（それは素晴らしい響きですね）

もちろん、他の何かでも全く同じように機能します。「爬虫類の宇宙」やすなわち「 Hornets の宇宙」について命題を立てることもできます。（それが住むのに魅力的な宇宙ではないでしょうか？）

これまでで分かったのは

-----------
                     |     |     |
                     |  1  |     |
                     |     |     |
                      -----------

は「一部の x と y」、つまり「一部の新しいものは美味しい」という意味であるということです。

あなたには追加の説明なしに、これは理解できると思います。

```
---------------------------
|     |     |
|     |  1  |
|     |     |
---------------------------

これは「一部のxはyである」という意味であり、つまり「一部の新しいものは好ましくない」を表しています。

今度はNo.5にグレーの駒を入れて、

```
---------------------------
|     |     |
|  0  |     |
|     |     |
---------------------------

を考えましょう。これはxとyの compartmentが空であることを示し、これを「すべてのxはyでない」と表現できます。つまり、「すべての新しいものは好ましくない」を表します。これがこのセクションの頭に書かれた三つの命題の二つ目です。

同様に、

```
---------------------------
|     |     |
|     |  0  |
|     |     |
---------------------------

は「すべてのxはyでない」という意味であり、つまり、「すべての新しいものは好ましくない」を表します。

このような状況はどう解釈すればよいか、興味深いですね？

```
---------------------------
|     |     |
|  1  |  1  |
|     |     |
---------------------------

これは二重の命題を表しています：つまり、「一部のxはyであり、また一部はyでない」、すなわち「一部の新しいものは好ましくあり、また一部は好ましくない」という意味です。

次に少し難しいかもしれません：

```
---------------------------
|     |     |
|  0  |  0  |
|     |     |
---------------------------

これは「すべてのxはyでなく、また何もyでない」、すなわち「すべての新しいものは好ましくなく、また何も好ましくない」という意味です。これにより、「新しいものは存在しない」という奇妙な結果が導かれます、つまり「ケーキは新しいものではない」となります。これは「好ましい」および「好ましくない」が「新しいケーキ」のクラスを完全に分割する「 Exhaustive 」分類であるためです：すなわち、これら二つはクラス全体を完全に網羅し、すべての存在する新しいケーキはどちらか一方に含まれる必要があります。

そして今度は、「すべてのケーキは新しいものでない」という命題の否定を考えましょう。これは「一部のケーキは新しいものである」、または文字を置き換えると「一部のケーキはxである」となります。どのように表現すればよいでしょうか？
```

この文は少し混乱させられるかもしれません。明らかに、新しいケーキがあることを知っているので、赤の玉をx半分の棚のどこかに入れなければなりません。しかし、左側の compartment には入れられません。なぜなら、それが美味しいかどうかは分かりません。また、右側の compartment も入れられません。なぜなら、それが不味いかどうかは分かりません。

では、どのようにすればよいのでしょうか？私はこの難局から抜け出す最善の方法として、xy-compartment と xy'-compartment の境界線上に赤の玉を置くことを提案します。これは、「1」を置くべき場所であると私がいつも表現するように、図で表します：

-----------
                     |     |     |
                     |    -1-    |
                     |     |     |
                      -----------

私たちの巧妙なアメリカの親戚は、どちらか一方の派閥に加入したいが、どの派閤に入るべきか決められない人の位置を表現する言葉を考案しました。たとえば、「民主党」と「共和党」の2つの派閥のいずれかに加入したいが、どちらに入れればよいか決断できないような人を「柵に座っている」と呼びます。この赤の玉が先ほど境界線上に置かれた位置はまさにそのようにです。彼は5番目に見えて好印象で、6番目も同様に見えて好印象ですが、どちらに入ればよいか決められません。そこで彼は柵の上で座りこみ、足を左右にぶら下げています！

次に、もっと難しい問題を与えます。これは何を意味しているのでしょうか？

-----------
                     |     |     |
                     |  1  |  0  |
                     |     |     |
                      -----------

これは明らかに二重の命題です。それは「一部のxはyである」というだけでなく、「すべてのxはNOT yである」ことも伝えています。したがって、結果として「すべてのxはyである」つまり「すべての新しいケーキは美味しい」となります。これはこのセクションの最初の3つの命題の最後です。

すると、全称命題

「すべての新しいケーキは美味しい」

は、二つの命題を組み合わせて形成されています。すなわち,

「一部の新しいケーキは美味しい」,
        そして「すべての新しいケーキは不味くない」.

同じようにして

```
---------------------------
|     |     |
|  0  |  1  |
|     |     |
---------------------------
```

これは「すべての x が y'である」という意味であり、

「新しいケーキは不味い」ということになります。

では、「君がくれたケーキは美味しい」という命題はどう解釈すればよいでしょうか？ 特称命題か普遍命題か？

もちろん、その答えは「特称」です。「一つのケーキだけは“何か”と呼ぶ価値がない」とすぐに答えるでしょう。

しかし、私の思いやり深い読者よ、それは「普遍的」である。覚えておくが、それらは少ないとはいえ（そして私はあなたがそれほど少ないとは思わないことを認めますが）、それが（正確には「それは」）君がくれたすべてのケーキだ！ したがって、「赤いものについては置いておいて」という条件を除くと、私のケーキの宇宙を二つのクラスに分けることにする。君がくれたケーキ（上段の棚に割り当てる）、そして君がくれなかったケーキ（下段に割り当てる）。すると下段はほぼ満杯で、上段はほとんど空っぽであることがわかる。そして、それぞれの半分に縦方向の区切りを入れて、美味しいケーキを左側、不味いケーキを右側に配置するという指示が与えられたとき、私はまず君がくれたすべてのケーキを集める（自分自身に言い聞かせるように、「優しい存在よ！ そんな親切なことをどう報いるべきか」と）。そしてそれを左側の compartiment に積み上げる。これには時間はかかりません！

もう一つの普遍命題を考えてみましょう。「バズリール・ベッカレッグは正直な男である」これは「現在考慮しているすべてのバズリール・ベッカレッグが正直な男である」という意味です。（私はその名前を作り出したと思っているでしょう？ しかし、そうではありません。それはコルnten のどこかの運送車に載っているのです。）

このような普遍命題（主語が単一のものである場合）は「個別」命題と呼ばれます。

では、「美味しいケーキ」という主語を用いて命題を考えましょう：つまり、棚の左側半分に注目し、そこにあるすべてのケーキが属性 y、すなわち「美味しい」を持っていることを固定しましょう。

---

もし、このようなマークが見つかったとしましょう：──|     |
                                                      |  1  |
何を意味しているのでしょうか？──|     |
                                    ────────────────
                                           |     |
                                           |     |
                                           |     |
                                    ────────────────

完全に横長の楕円形について説明した後で、縦長の楕円形については長い時間を費やす必要はないと願っています。あなた自身で理解できると思いますが、これは「一部の y が x である」という意味であり、

「一部の美味しいケーキが新しい」

という意味です。

しかし、「それでは以前にも同じ状況がありました」と言うでしょう。「赤い玉を第5番目に置きましたし、それは『一部の新しいケーキが美味しい』という意味だと教えたのです。そして今度は『一部の美味しいケーキが新しい』と意味していると言いますか？両方とも意味するのでしょうか？」

この質問は非常に深いもので、読者様には大変素晴らしいご質問をなさったと思います！それは確かに両方とも意味します。もし x（つまり、「新しいケーキ」）を選んで主題として取り扱い、第5番目を横長の楕円形の一部と見做すなら、「一部の x が y である」という解釈ができる、すなわち「一部の新しいケーキが美味しい」と読むことができます。しかし、もし y（つまり、「美味しいケーキ」）を選んで主題として取り扱い、第5番目を縦長の楕円形の一部と見做すなら、「一部の y が x である」という解釈ができる、すなわち「一部の美味しいケーキが新しい」と読むことができます。これらは単に同じ真実を異なる方法で表現するだけです。

言葉を使わず、他のこの縦長の楕円形がどのようにマークされるか、そしてそれぞれの場合における意味を記しておきます。これらの表と横長の楕円形のさまざまな状況を比較することで、私はあなたが明確に理解できると思います。

自分自身でこのテーブルをテストする良い方法は、最初一つの列目を隠し、次に別の列目を隠しながら「ダッキング」することです。また、下半分と右半分の棚についてそれぞれ別のテーブルを作成することもお勧めします。

そして今では、私たちは小さな図形について十分な説明を終えていると思いますので、大きな図形に進むことにしましょう。

---

記号．意味．
-------------------------
|-----|--------------------------------------|
| --- | ------------------------------------ |
| --- | 一部の y は x' である；               |
| --- | 即ち、一部の美しいものは新しいものでない．|
|-----|--------------------------------------|
| --- | ------------------------------------ |
| 1   | ------------------------------------ |
| --- | ------------------------------------ |
|-----|--------------------------------------|
| --- | ------------------------------------ |
| --- | 全ての y は x である；                 |
| --- | 即ち、全ての美しいものは新しいものでない．|
| --- | [これは「全ての新しいものは美しいものではない」という表現の別の形です]|
|-----|--------------------------------------|
| --- | ------------------------------------ |
| --- | ------------------------------------ |
| 1   | ------------------------------------ |
| --- | ------------------------------------ |
|-----|--------------------------------------|
| --- | ------------------------------------ |
| --- | 全ての y は x' である；                |
| --- | 即ち、全ての美しいものは新しいものでない．|
| --- | ------------------------------------ |
| 0   | ------------------------------------ |
| --- | ------------------------------------ |
|-----|--------------------------------------|
| --- | ------------------------------------ |
| --- | 一部の y は x であり、一部は x'；       |
| --- | 即ち、一部の美しいものは新しいもので、一部は新しいものではない．|
| 1   | ------------------------------------ |
| --- | ------------------------------------ |
|-----|--------------------------------------|
| --- | ------------------------------------ |
| --- | 全ての y は x である、そして全てが x'；|
| --- | 即ち、美しいものは存在しない．          |
| --- | 即ち、ケーキは美しいものではない．     |
| 0   | ------------------------------------ |
| --- | ------------------------------------ |
|-----|--------------------------------------|
| --- | ------------------------------------ |
| --- | 全ての y は x；                        |
| --- | 即ち、全ての美しいものは新しいものである．|
| --- | ------------------------------------ |
| 0   | ------------------------------------ |
| --- | ------------------------------------ |
|-----|--------------------------------------|
| --- | ------------------------------------ |
| --- | 全ての y は x'；                       |
| --- | 即ち、全ての美しいものは新しいものではない．|
| --- | ------------------------------------ |
| 1   | ------------------------------------ |
| --- | ------------------------------------ |
|-----|--------------------------------------|

この部分は、最後の部分と同じように分割された棚とみなすことができますが、Attribute mのためにも二分されています。mを「栄養価のある」という意味に与えましょう。そして、すべての「栄養価のある」ケーキが中央の正方形内に置かれ、すべての「栄養価のない」ケーキはその外側、つまり四つの不規則な形状の外側 compartiment に置かれるという仮定を立てましょう。

これにより、各 compartiment 内のケーキには三つの属性が付与されることになります。また、二つの属性を表す文字は以前の図形の場合と同様に compartiment の辺缘に記載されますが、ここでは角に記載されます。（注意していただきたいのは、m' は四つの外側の角全てに書かれることとします。）これにより、各 compartiment 内の物体が持つ三つの属性を一目で確認することができます。

たとえば、No.12 を見てみましょう。その角には x, y', m が記載されていますから、この compartiment 内のケーキ（もし存在する場合）は「新しい」「味が良くない」そして「栄養価のある」という三属性を持つことがわかります。また、No.16 を見てみましょう。その角には x', y', m' が記載されていますから、この compartiment 内のケーキ（もし存在する場合）は「新しいでない」「味が良くない」そして「栄養価のない」という三属性を持つことがわかります。（実に魅力的ではないケーキたちです！）

この図形上に表される x と y, x と m, y と m を含むすべての命題を一つずつ確認することは非常に時間がかかります（合計で96個ありますから、ご容赦いただければと思います）ので、私は二つや三つの例だけを示すことにします。自分自身で多くのものを計算してみることをお勧めします。

上半分だけを取り出して、「新しいケーキは存在しない」と表現する方法を考えましょう。これは文字を使って表すと「x が m でない」ことになります。これにより、棚の上半分に属するケーキのうち、中央の正方形内に存在することはなく、No.11 と No.12 の二つの compartiment は空であることがわかります。そしてこれは、以下の図示によって表現されます：$$\begin{array}{|c|c|}
\hline
 & m' \\
\hline
x' &  \\
\hline
 & y' \\
\hline
\end{array}$$

---

-------------------
                  |         |         |
                  |    \_\_\_\_\_|\_\_\_\_\_    |
                  |   |     |     |   |
                  |   |  0  |  0  |   |
                  |   |     |     |   |
                   -------------------

そして、矛盾する命題「 SOME x は m である」をどのように表現するべきか。これは以前に考慮した難点である。私は最も良い方法として、赤い駒を第11番と第12番の境界線上に置くことと考えている。これは一方の compartiment が「占領されている」ということを意味し、現在はどちらが占領されているかは不明である。これを以下のように表現する。

-------------------
                  |         |         |
                  |    \_\_\_\_\_|\_\_\_\_\_    |
                  |   |     |     |   |
                  |   | -1- |     |   |
                  |   |     |     |   |
                   -------------------

次に、「すべての x は m である」という命題を表現する方法について考えよう。これは、二つの命題から成る。

" SOME x は m である,"
             そして " NO x は m' である。"

まず否定的な部分を表現しよう。これは上半分の棚に属するケーキが中央正方形の外側にあることがないことを示す。つまり、第9番と第10番の compartiment は空である。これは以下のようになる。

-------------------
                  | 0       |       0 |
                  |    \_\_\_\_\_|\_\_\_\_\_    |
                  |   |     |     |   |
                  |   |     |     |   |
                  |   |     |     |   |
                   -------------------

しかし、「 SOME x は m である」を表現する方法についてまだ考えなければならない。これは第11番と第12番の oblong 内にいくつかのケーキがあることを示すので、前例のように、第11番と第12番の境界線上に赤い駒を置き、結果は以下のようになる。

-------------------
                  | 0       |       0 |
                  |    \_\_\_\_\_|\_\_\_\_\_    |
                  |   |     |     |   |
                  |   | -1- |     |   |
                  |   |     |     |   |
                   -------------------

最後に、一つまたは二つの解釈を試してみよう。

---

この図について、xとyに関してどのように解釈すべきか考えてみましょう。

-------------------
                  |         |       0 |
                  |    \_\_\_\_\_|\_\_\_\_\_    |
                  |   |     |     |   |
                  |   |  1  |  0  |   |
                  |   |     |     |   |
                   -------------------

この図は、xy'の正方形に関して、両方の compartiment が「空」であることを示しています。一方で、xyの正方形に関しては、「占領されている」という情報を得ます。確かに、その compartment は一つだけですが、他の compartment の状態に関係なく、正方形内に何かがあることは確定します。

したがって、これらのマークを小さな図形に移動し、m分割を取り除くと、

-----------
                     |     |     |
                     |  1  |  0  |
                     |     |     |
                      -----------

というマークが可能です。これは「すべての x は y」を意味します。

もし与えられた長方形が以下のようにマークされていた場合でも、結果は同じです：

-------------------
                  | 1       |       0 |
                  |    \_\_\_\_\_|\_\_\_\_\_    |
                  |   |     |     |   |
                  |   |     |  0  |   |
                  |   |     |     |   |
                   -------------------

再度：xとyに関して、このマークはどう解釈すべきか考えてみましょう。

-------------------
                  | 0       |       1 |
                  |    \_\_\_\_\_|\_\_\_\_\_    |
                  |   |     |     |   |
                  |   |     |     |   |
                  |   |     |     |   |
                   -------------------

この図は、xyの正方形に関して、その compartment の一つが「空」であることを示しています。しかし、他の compartment にはマークがないため、この情報だけでは何も結論を出せません。もし他の compartment も「空」であれば、正方形全体も「空」になりますし、もし「占領されている」であれば、正方形は「占領されている」ことになります。したがって、どちらの状況であるかわからないため、この正方形に関して何とも言えません。

他方の正方形、すなわち \(xy'\)-正方形については、（先ほどの例と同様に）「占領されている」と理解している。

したがって、我々のマークを小さな図形に移動させると、以下のような状況になる：

\[
\begin{array}{|c|c|c|}
\hline
 &  &  \\
\hline
 & 1 &  \\
\hline
 &  &  \\
\hline
\end{array}
\]

これは「ある \(x\) が \(y'\)」という意味である。

これらの原理はすべての他の長方形にも適用できる。例えば、「すべての \(y'\) が \(m'\)」を表現するには、属性 \(y'\) を持つ右側縦長方形（以下に示すようなもの）を次のようにマークする：

\[
\begin{array}{|c|}
\hline
\text{} \\
\hline
\text{-1} \\
\hline
0 \\
\hline
\end{array}
\]

また、もし下半分の棚が以下のようにマークされていて、\(x\) と \(y\) の関係について解釈を求められた場合、

\[
\begin{array}{|c|c|c|c|}
\hline
 &  &  & \text{} \\
\hline
 & 0 &  & \text{} \\
\hline
 &  &  & \text{} \\
\hline
1 & \text{-} & \text{0} & \text{} \\
\hline
\end{array}
\]

以下のように小さな図形に移動させ、解釈する：

\[
\begin{array}{|c|c|}
\hline
 &  \\
\hline
1 & 0 \\
\hline
 &  \\
\hline
\end{array}
\]

これは「すべての \(x'\) が \(y\)」という意味である。

命題に関するもう2つの注意点を述べる必要がある。

一つ目は、すべての「ある」または「すべての」と始まる命題において、実際的な存在性が主張されていることである。例えば、「すべての miser（avaricious人）が selfish（selfishな人）である」という命題を述べた場合、これは miser が実際に存在することを意味している。もしこの存在性を否定し、むしろ miserliness（avariciousness）が necessarily involves selfishness（selfishであることを必要とすること）という法則だけを述べたい場合は、「すべての miser は unselfish（selfishでない人）ではない」という命題を使うべきであり、これは miser の存在性を主張していないが、もし存在した場合でも必ず selfish であるということだけを述べている。

他の場合では、命題が「一部」または「全くない」と始まり、さらに二つの属性を超える属性を含むとき、これらの属性は任意に並び替えられ、一つの項から別の項へ移動することができます。「自由に」としてです。例えば、「一部abcはdefである」という命題は、「一部bfはacdeである」と再配置され、それぞれが「一部のものはabcdefである」と等価であることを示します。また、「全くない賢明な老人は冒険的な賭博者である」は「全くない冒険的な賭博者の老人は賢明で冒険的である」と再配置され、それぞれが「全くない男性は賢明な老人で冒険的で不注意な賭博者である」と等価であることを示します。

2. シュロジズム

今、我々の事象の宇宙を三つの異なる属性に関して三分するものとしましょう。これらの三つの属性から、三つの異なる組み合わせを作ることができます（たとえば、a, b, cが与えられた場合、ab, ac, bcという三つの組み合わせを作ることができます）。また、二つの命題が与えられ、これらはこの三つの組み合わせのうちの二つを含むものとします。そして、それらから第三の命題（その他の組み合わせを含むもの）を証明することができます。（たとえば、m, x, yという事象の宇宙を分ける場合、与えられた二つの命題が「全くないmはx'である」と「すべてのm'はyである」であり、これらはmxとmyという二つの組み合わせを含むものとします。それらから第三の命題（xとyを含むもの）を証明することができます）

このような場合、与えられた命題を「前提」と呼び、第三のものを「結論」とし、全体を「シュロジズム」と呼ぶことがあります。

明らかに、一つの属性は両方の前提で出現する必要があります。または、一つの属性が一方の前提で現れ、その対偶が他方の前提で現れる場合もあります。

最初の場合（たとえば、前提が「一部mはxである」と「全くないmはy'である」の場合）、両方に現れる項を「中間項」と呼びます。これは他の二つの項を結ぶ役割を果たすからです。

第二の場合（たとえば、前提が「全くないmはx'である」と「すべてのm'はyである」の場合）、これらの対偶属性を含む二つの項を「中間項」と呼びます。

したがって、最初の場合では、「m-事象」のクラスが中間項であり、第二の場合では、「m-事象」のクラスと「m'-事象」のクラスが中間項となります。

中間項または中間項に現れる属性は結論で消え、その存在を「排除」と呼びます。これは字義通りに「追い出される」という意味です。

現在、二つの前提から結論を導こうとしましょう—

「いくつかの新しいケーキは不健康である；
  すべての丁寧なケーキは不健康ではない。」

これらの前提を石で表現するには、ケーキを新しさ、丁寧さ、そして不健康さという三点に分類する必要があります。そのためには、大きな図形を使用し、x を「新しい」と、y を「丁寧な」と、m を「不健康な」で表す必要があります。（中央の正方形内のすべてのものは属性 m とされ、その外側は m'（つまり「非-m」）とされます。）

中項や中間項が示す属性を意味する規則に従って、m を使用することをお勧めします。（私は m をシンボルとして選んだのは、「中間」が「m」で始まるからです。）

二つの前提を表現する際は、否定的な前提（「ない」と始まるもの）から始めましょう。灰色の石は常に確実に配置でき、その後赤い石の位置を固定する手助けになります。これらの石の位置は時として少し不確定であるため、最も歓迎される場所を特定するために有用です。

「すべての丁寧なケーキは不健康ではない（ケーキ）」と表現します。つまり、「y-ケーキはm'-ケーキではない」という意味です。これは、棚の y の半分に属するケーキのうち、中央の正方形の外側にある m' の compartment には何も含まれていないことを示しています。したがって、 compartment No. 9 と No. 15 はどちらも「空」であり、それぞれに灰色の石を配置します。

以下のように配置します：

```
-----------
|0    |     |
|   --|--   |
|  |  |  |  |
|--|-----|--|
|  |  |  |  |
|   --|--   |
|0    |     |
-----------
```

現在、残りの前提を表す必要があります。つまり、「新しいケーキは不健康である（ケーキ）」というものです。これは「一部のx-ケーキはm'-(ケーキ)である」という意味です。この文は、棚のx-半分にある一部のケーキがm'- compartimentに存在することを示しています。したがって、9番と10番の2つの compartimentのうちの一つは「占領」されていることがわかります：しかし、どの compartmentに赤い駒を置くべきかについては明示されていないため、通常のルールではその境界線上に置くのが一般的です。しかし、この場合、他の前提が私たちを助けており、9番が空であると宣言したので、赤い駒は選択肢を持っていませんし、10番に配置される必要があります：

```
 -----------
 |0    |    1|
 |--|--   |
 |  |  |  |
 |--|-----|--|
 |  |  |  |
 |   --|--   |
 |0    |     |
 -----------
```

次に、この情報を利用して小さな図形にxとyのみを含む命題を得るためにどの駒を配置するかを見ていきましょう。4つの compartimentを一つずつ見てみます。

まず、5番の compartmentです。これについて知っているのはその外側が空であることです：しかし、内側については何も知りません。したがって、正方形は空である可能性もありますし、何かがあるかもしれません。誰もそれを知る方法はありませんので、この正方形には何の駒も置けません。

次に6番の compartmentです。ここでは少し状況が改善されています。外側に赤い駒があることを知っているため、内部が空か占められているかは分かりません：しかし、それがどう関係するのでしょうか？ 一方の角にある単一のケーキは十分な理由となり、「この正方形は占められています」と言えるし、それを赤い駒でマークすることができます。

7番の compartmentについては5番と同様に状況が同じです—部分的に「空」であることがわかりますが、その他の部分が空か占められているかは分かりません：したがって、この正方形には何も描けません。

そして8番の compartmentについては全く情報がないです。

結果は以下のようになります：

```
 -------
 |   | 1 |
 |--|---|
 |   |   |
  -------
```

私たちの「結論」は、xy'方格に赤い駒があるという非常に少ない情報から導き出す必要があります。したがって、私たちの結論は「一部のxはy'である」となり、「一部の新しいケーキは不味い（ケーキ）である」となります：または、もしあなたがy'を主語として選ぶのであれば、「一部の不味いケーキは新しい（ケーキ）である」ですが、前者の方が見栄えが良いです。

次に、以下の論理全体を書き出しましょう。符号\&there4[*]で「従って」と置き換え、各命題の末尾に「ケーキ」を省略して簡潔にします。

[*][ブレットからの注記：\&there4というシンボルの使用は比較的任意な選択です。一般的なフォントでは、「従って」のシンボル（三角形内に並んだ3つの点）が正確には表示できませんが、HTMLでこのシンボルを適切に認識することができます。この文書がブラウザで読まれる場合、そのシンボルは正しく認識されるでしょう。これは貧者の言い訳です.]

「一部の新しいケーキは不味い；
すべての良いケーキは不味くない
従って、一部の新しいケーキは不味くない。」

あなたは、初めての「三段論法」を成功裏に解きほぐしました。私からお祝いの言葉を申し上げますし、また、これは長いそして栄光に満ちた類似した勝利の一連の最初であることを希望します！

もう一つの三段論法を解きほぐしてみましょう—これは最後よりもやや難しいものですが—その後、私はあなたが自分自身でゲームを楽しむことができると考えます（より良いのは、友人を見つけ、その友人がこのスポーツに参加する意欲と能力がある場合です）。

では、次の前提から何を導き出すことができるでしょうか？

「すべてのドラゴンは不気味；
すべてのスコットランド人は気味が良い。」

注意してください、私はこれらの前提が事実であることを保証していません。まず第一に、私はドラゴンを見たことがありません：そして第二に、私たちが論理学者として、これらの前提が真偽に関わらず何らかの影響を持つこと自体は微塵も関係ありません：私たちはただ、これらが結論を論理性で導くかどうかを証明する必要があります。つまり、もしこれらが真であれば、その結論も真であることを示す必要があります。

見よ、現在「ケーキ」を手放さなければならず、それによって棚は私たちにとって無用になる。そこで、「宇宙」として何かのクラスを選ばなければならない：ドラゴンやスコットランド人を含むものとして「動物」を選びましょう。そして、「賢明」という属性が明らかに「中間項」に属するものであるから、m を「賢明」、x を「ドラゴン」、y を「スコットランド人」とします。したがって、私たちの二つの前提は、

「すべてのドラゴン動物は不賢明（動物）；
すべてのスコットランド人動物は賢明（動物）。」

となります。

これらの前提を文字で表すと、

「すべての x は m'；
すべての y は m。」

となります。

最初の前提は、あなたがすでに知っているように、二つの部分から成ります：

「一部の x は m'」
そして、「すべての x は m」。

第二前提も同様に二つの部分から成ります：

「一部の y は m」
そして、「すべての y は m'」。

まず否定的な部分から始めましょう。

そこで、大きな図形上では最初に「すべての x は m」をマークし、次に「すべての y は m'」をマークします。もう説明不要でしょうが、それぞれ独立して、

```
-----------
|     |     |
|   --|--   |
|  |0 | 0|  |
|--|--|--|--|
|  |  |  |  |
|   --|--   |
|     |     |
 ----------- 
```

と

```
-----------
|0    |     |
|   --|--   |
|  |0 | 0|  |
|--|--|--|--|
|  |  |  |  |
|   --|--   |
|0    |     |
 ----------- 
```

となります。これらの二つを組み合わせると、

```
-----------
|0    |     |
|   --|--   |
|  |0 | 0|  |
|--|--|--|--|
|  |  |  |  |
|   --|--   |
|0    |     |
 ----------- 
```

となります。

次に、二つの肯定的な部分、「一部の x は m'」と「一部の y は m」をマークします。

x が xm' を持つ唯一の二つの compartment は No.9 と No.10 です。No.9 は既に「空」でマークされているため、赤い駒は No.10 に入れます。

同様に、y が ym を持つ唯一の二つの compartment は No.11 と No.13 です。No.11 は既に「空」でマークされているため、赤い駒は No.13 に入れなければなりません。

最終的な結果は以下の通りです

\[
\begin{array}{|c|c|}
\hline
0 & 1 \\
\hline
-- & -- \\
\hline
|0 | 0| \\
\hline
-- & -- & -- \\
\hline
|1 | \quad | \\
\hline
0 & \quad \\
\hline
\end{array}
\]

そして、この情報のうちどれくらいが有用に小さい図形に移すことができるでしょうか？

まず、その四つの部分を一つずつ見てみましょう。

No. 5については？これは完全に「空」です。（灰色の駒でマークします）

No. 6については？これは「満たされている」状態です。（赤色の駒でマークします）

No. 7については？同じく「満たされている」状態です。

No. 8については？情報はありません。

現在、小さい図形はかなり詳しくマークされています：

\[
\begin{array}{|c|c|}
\hline
0 & 1 \\
\hline
--- & --- \\
\hline
1 & \quad \\
\hline
\end{array}
\]

そして、この情報を基にどのような結論を読み取ることができるでしょうか？実際には、このような豊富な情報は一つの命題では表現できないので、二つの命題が必要です。

まず、xを主語として「すべての x が y'」とします。つまり,

「すべてのドラゴンはスコットランド人ではない」:

次に、yを主語として「すべての y が x'」とします。つまり,

「すべてのスコットランド人はドラゴンではない」.

これら二つの前提と結論を一緒に書き出しましょう。

\[
\text{"すべてのドラゴンは不気味である; }
\quad \text{すべてのスコットランド人は明瞭である.}
\quad \therefore \text{すべてのドラゴンはスコットランド人ではない;}
\quad \text{すべてのスコットランド人はドラゴンではない."}
\]

最後に、論理に関する文献では、存在する何かが全くないという前提で議論される場合があります。つまり、「ある x が y」は「属性 x, y は互換的であり、何か一つには両方を同時に持つことができる」と解釈され、「すべての x が y」は「属性 x, y は互換的ではないので、何も両方を持てない」と解釈されます。

このような文献では、命題の意味は私たちの「論理ゲーム」におけるものとは大きく異なります。その違いを正確に理解することが重要です。

まず「some x are y」について考えます。ここでは、「are」という言葉を「実際の事実として存在すること」と理解します。もちろん、これはx-事物が存在することを意味します。しかし、他の論文の著者は、「are」という言葉を「可能である」と解釈し、これには何らxが存在することを含むものではありません。したがって、彼らは私たちよりも少ない意味を持ちます：私たちの意味は彼らの意味を含みますが（もちろん、「some x ARE y」は「some x CAN BE y」を含む）、彼らの意味は私たちの意味を含まない。

次に、「no x are y」について考えます。ここでは、「are」という言葉を「実際の事実として存在すること」と理解します。これはxがyであることを意味するものではありません。しかし、他の論者はこの命題を解釈し、何もyでないだけでなく、何らかの形でyになることも不可能であると主張します。したがって、彼らは私たちよりも多くの意味を持ちます：彼らの意味は私たちの意味を含みますが（もちろん、「no x CAN be y」は「no x ARE y」を含む）、私たちの意味は彼らの意味を含まない。

例えば、「some Welsh hippopotami are heavy」は、他の論者の解釈では真ですが（「Welsh」と「heavy」がヒッポポタマスに対して互換的であるため）、私たちのゲームでは偽です（実際にはウェールズのヒッポポタマスが存在しないから）。

さらに、「all x are y」について考えます。これは二つの部分命題「some x are y」と「no x are y'」で構成されます。もちろん、最初の部分では私たちよりも少ない意味を持ちますが、第二の部分では私たちよりも多くの意味を持ちます。しかし、これらの操作は相互に平衡することはなく、あなたが一つの煙突を倒したことを慰めるために追加のドアステップを与えるのとは似て非なるものです。

もしこのような三段論法に遭遇した場合、私が与えたシステムを使って計算することは非常に簡単です。ただ「ある」を「可能である」と意味するようにすれば、スムーズに進むことでしょう。つまり、「一部の x は y である」という文は「一部の x は y の可能性がある」となり、「属性 x, y は互換的である」となります。そして、「すべての x は y でない」という文は「すべての x は y の可能性がない」となり、「属性 x, y は不互換的である」となります。もちろん、「すべての x は y である」という文は「一部の x は y の可能性があるし、一部の x は y' の可能性がない」となり、「属性 x, y は互換的であり、属性 x, y' は不互換的である」となります。このシステムのために図形を使用する際には、赤い駒を「この compartiment に何かがある可能性がある」と意味し、灰色の駒を「この compartiment に何も存在しない可能性がない」と理解してください。

3. 誤謬

そしてあなたは本当にそう思っているのですか？実生活における論理の主な目的は、作業可能な前提から結論を導き出し、他人が導いた結論が正しいことを確認することではないでしょうか？それが本当であれば素晴らしいです！社会はパニックや他の幻覚に非常に少なくなっており、特に政治生活は全く異なるものになっているはずです。しかし、世界中で散見される議論のうち、正確なものはわずかであることが恐れられています。あなたが新聞や雑誌を読むときに出くわす作業可能な一対の前提（つまり、論理的な結論を導き出す一対）に対して、おそらく5つは全く結論を導けないものでしょう。そして、前提が作業可能である場合でも、作者が正確な結論を導いた例に対して1つしかない一方で、不正確な結論を導いた例は10個以上あることが多くなるのです。

最初の場合には、「前提が誤謬的である」と言えますし、後者の場合は「結論が誤謬的である」と言えます。

このゲームによって学んだ論理的なスキルの主な用途は、これらの2種類の誤謬を検出することにあります。

初めの一種の論理誤謬、「虚偽の前提」は、大きな図形にマークを付けた後、それを小さな図形に移す試みで発見します。四つの compartiment 一つずつを確認し、それぞれについて「ここにはどのようなマークを置くことができるか？」と尋ねます。そして、これら全てにおいて答えは「情報なし！」となり、結論が存在しないことを示します。例えば、

```
すべての兵士は勇敢である；
一部のイギリス人は勇敢である。
従って、一部のイギリス人は兵士である。
```

これは典型的な三段論法に見えて、経験の浅い論理学者には容易に誤解されるかもしれません。しかし、あなたはそのような罠にはまりません！前提を書き出し、冷静に「虚偽の前提！」と述べます。作者が主張した結論について尋ねる必要はありません——どんなものであれ、それは間違いであることが分かっているからです。賢明な母親が言っていたように、「マリー、保育室に行き、赤ちゃんが何をしているか見てきて、彼にそれをしないために伝えてくるんだ」というような態度で対処します。

もう一つの論理誤謬、「虚偽の結論」は、両方の図形をマークし、正しい結論を読み取り、作者が導き出した結論と比較した後で発見されます。

ただし注意してください。正確な結論とは異なるだけでは「虚偽の結論」と言わないでください：それは正確な結論の一部である可能性があり、その部分に関しては完全に正しいかもしれません。この場合、「欠陥のある結論！」と同情的な笑顔を浮かべながら述べます。例えば、

```
すべての無私心の人々は大方である；
一切のavaricious人は大方ではない。
従って、一切のavaricious人は無私心の人々ではない。
```

という三段論法に遭遇した場合、前提は以下のように文字で表現できます：

```
すべてのx'はmであり；
一切のyはmではない。
```

ここでは正しい結論は「すべてのx'はy'である」（つまり、「すべての無自覚の人々は非貪婪者である」ということ）であり、著者は結論として「すべてのyはx'でない」（これは「すべてのx'はyではない」と同じで、「すべてのx'はyである」の一部である）を導き出している。ここでは単に「不完全な結論！」と言えばよいだろう。同様のことが、菓子店で、小さな男の子が二便士を置いて、勝利気分で一つのパンケーキを持って帰っていった場合にも起こる。あなたは悲しげに頭を振って、「不完全な結論！」と述べるだろう。そしておそらくカウンターの若い女性に尋ねてみるかもしれない。「この小さな男の子が支払いを済ませた後で残したパンケーキを君に食べさせてくれないか？」と。そしておそらく彼女は「そうではない」と答えるだろう。

しかし、前述の例で著者が結論として「すべての貪婪者は無自覚である」（つまり、「すべてのyはxである」ということ）を導き出した場合、これは彼の正当な権利を超えたものであり（これによりyの存在が主張され、前提には含まれていない）、あなたは適切に「論理的誤りの結論！」と言えるだろう。

読書中に他の論理に関する文献を読むと、さまざまな種類の（そう呼ばれる）「論理的誤り」に遭遇するかもしれない。しかしこれらが必ずしも常にそのようなものではないことに注意せよ。例えば、「誠実な人間は詐欺をしない；不誠実な人間は信頼性がない」という前提を与えられたロジシャンに対して、結論は何であるかと尋ねた場合、彼はおそらく「何もありません！ あなたの前提は二つの異なる規則に違反しており、どれほど論理的に欠陥があるかを示しています！」と言うことだろう。そしてあなたが勇気を持って「詐欺をする人間は信頼性がない」という結論であると述べた場合、私は懸念する。あなたの論理的な友人は急いで去るかもしれない—激怒しているかもしれませんし、ただ嗤っているだけかもしれません。いずれにせよ結果は不快なものとなるだろう。私にはそのような実験を試みないことを強く推奨する。

「ではなぜそうなるのか？」とあなたは言うだろう。「あなたが私たちにすべてのこれらのロジシャンが間違っていると言いたいのですか？」全くそんなことはない、親愛なる読者！ 彼らから見れば完全に正しい。しかし彼らの体系には、全ての可能性のある三段論法の形を含むわけではありません。

彼らは否定接頭辞で始まる属性に対して一種の神経質な恐怖心を持っている。例えば、「すべての非xはyである」や「すべてのxは非yではない」というような命題は、彼らの体系からは完全に外れている。そして、この恐怖から一時的に有用な多くの形を排除した結果、彼らは少数の許容される形式に対応する規則を作り出したが、これらは全ての可能な形に対して全く役立たない。

我々は読者よ、争う必要はありません！世界には私たちの体系も十分に余裕があります。私たちはより広い体系を静かに受け入れましょう。もし彼らがこれらの有用な形式を見過ごし、「これらは完全に三段論法ではない！」と言ったとしても、それに対して我々は立ち去るしかありません。彼ら自身の運命に向かって突進させることでしょう。あなたが自分の運命に対して急いで行動することは非常に危険です。ポテトベッドやストロベリーベッドに急いで入ることは大きな被害を被る可能性はありませんし、バルコニーにも急いで入ること（建設中でない限り）は生き延びることが可能です。しかし、一度運命に対して急いで行動すると、その結果を受け入れるしかありません。

\newpage

\chapter{問いかけの数々（Cross Questions）}
\label{ch:ii}

\begin{center}
\large\textit{大小ダイアグラムへの命題配置・記号解釈と演習問題体系}
\end{center}
\vspace{4mm}

第二章

交錯問題

「荒野の男は私に尋ねた。
『海の中に何個のストロベリーが生えているか？』」

1. 初等的。

1. 何かを「属性」とは何でしょうか？例を挙げて説明してください。

2. 「は」または「です」を二つの名詞間に置くことが意味のある理由はいつでしょうか？例を挙げて説明してください。

3. 「は」または「です」を置くことが意味がない理由はいつでしょうか？例を挙げて説明してください。

4. 「は」または「です」を置くことが意味がない場合、最も単純な合意は何ですか？

5. 何が「命題」、「項」、「主語」、「述語」でしょうか？それぞれの例を挙げて説明してください。

6. 「個別」と「普遍」の命題とは何でしょうか？それぞれの例を挙げて説明してください。

7. 小さい図形を見たときに、各 compartiment に属する属性を知るための規則は何ですか？

8. ロジックにおいて、「some」とは何を意味しますか？[参照：第55ページ，第6ページ]

9. このゲームにおける「宇宙」はどのような意味で使用されますか？

10. 「双対」の命題とは何でしょうか？例を挙げて説明してください。

11. 何か事物のクラスが「完全に」分割されるとはいつでしょうか？例を挙げて説明してください。

12. 「柵に座っている」という表現はどのように解釈されますか？

13. 「すべて x は y である」という命題を構成する二つの部分的な命題は何ですか？

14. 「個別」の命題とは何でしょうか？例を挙げて説明してください。

15. このゲームにおいて、どのような種類の命題がその主語の存在を含意しますか？

16. 命題に二つ以上の属性が含まれている場合、これらの属性は一部の場合には再配置され、一つの項から他の項へ移動することができます。どの場合にこれを行うことができますか？例を挙げて説明してください。

\_\_\_\_\_\_\_\_\_\_

次の各命題を二つの部分的な命題に分解してください：

17. すべての虎は凶暴である。

18. すべての茹で卵は不健康である。

19. 私は幸せである。

20. ジョンは家にいない。

\_\_\_\_\_\_\_\_\_\_

[参照：第56ページ，第7ページ]

21. 大きな図形を見たときに、各 compartiment 内に含まれる事物の属性を知るための規則は何ですか？

22. 何が「前提」、「結論」、「三段論法」でしょうか？それぞれの例を挙げて説明してください。

23. 「中間項」という表現と「中間項」について説明してください。

24. 大きな図形に一対の一連の前提をマークするとき、否定的な命題を肯定的な命題よりも前にマークすることが最も良い理由は何ですか？

25. ロジシャンとして、前提が真偽であるかどうかは私たちにとって何の関心事でもないという理由は何ですか？

26. 「一部のxがyである」とは「属性xとyは互換的である」と解釈し、「一部のxがyでない」とは「属性xとyは不互換的である」とはどういう論理を表しているか？

27. 二つの種類の「誤謬」は何であるか？

28. 「誤謬的な前提」を見つける方法は何か？

29. 「誤謬的な結論」を見つける方法は何か？

30. 提出された結論が正しい結論とは異なる場合でも、公平には「誤謬的」と呼べないことがある。そのような状況はどのようなときか？また、このような結論には何という名前を付けるべきか？

[参照：第57-59ページ]

---

3. 小図形の半分。

表現すべき命題。

```
----------------
|     |     |
|     x     |
|     |     |
--y-----y'--
```

1. 一部のxはyでない。
2. 全てのxはyでない。
3. 一部のxはyであり、また一部はyではない。
4. xは存在しない。
5. 一部のxが存在する。
6. 一切のxはyでない。
7. 一部のxはyでないし、また一部のxが存在する。

\_\_\_\_\_\_\_\_\_\_

xを「判事」、yを「正直」とする。

8. 判事は正直ではない。
9. 一部の判事が不正直である。
10. 全ての判事が正直である。

\_\_\_\_\_\_\_\_\_\_

xを「梅」、yを「栄養価のあるもの」とする。

11. 一部の梅は栄養価がある。
12. 栄養価のある梅は存在しない。
13. 梅は一部は栄養価があり、また一部はそうではない。
14. 全ての梅は栄養価がない。

[参照：第59-60ページ]

```
-------------
|     |     
|     x
|     |
|--y--|
|     |
|     x'
|     |
-------------
```

\_\_\_\_\_\_\_\_\_\_

yを「怠け者でない学生」、xを「成功した」とする。

15. 一切の怠け者でない学生は失敗していない。
16. 全ての怠け者でない学生は成功している。
17. 学生は一切怠け者ではない。
18. 一部の怠け者でない学生は失敗している。
19. 一部の学生が怠け者である。

[参照：第60-1ページ]

---

-------            -------
                 |   |   |          |   |   |
             1.  |   | 0 |      2.  | 0 | 0 |
                 |   |   |          |   |   |
                  -------            -------

-------            -------
                 |   |   |          |   |   |
             3.  |   -   |      4.  | 0 | 1 |
                 |   |   |          |   |   |
                  -------            -------

\_\_\_\_\_\_\_\_\_\_

x="好問の謎"; y="難問"として、

-------            -------
                 |   |   |          |   |   |
             5.  | 1 |   |      6.  | 1 | 0 |
                 |   |   |          |   |   |
                  -------            -------

-------            -------
                 |   |   |          |   |   |
             7.  | 0 | 0 |      8.  | 0 |   |
                 |   |   |          |   |   |
                  -------            -------

\_\_\_\_\_\_\_\_\_\_

[参照：p.61, 2]

x="蟹"; y="自己中心的"として、

-------            -------
                 |   |   |          |   |   |
             9.  |   | 1 |     10.  | 0 |   |
                 |   |   |          |   |   |
                  -------            -------

-------            -------
            11.  | 0 | 1 |     12.  | 1 | 1 |
                 |   |   |          |   |   |
                  -------            -------

\_\_\_\_\_\_\_\_\_\_

-----
                         |     |
                         x     |
                         |     |
                         |--y'-|
                         |     |
                         x'    |
                         |     |
                          -----

y="健康な人々"; x="幸せ"として、

---          ---          ---          ---
           | 0 |        |   |        | 1 |        | 0 |
      13.  |---|   14.  |-1-|   15.  |---|   16.  |---|
           | 1 |        |   |        | 1 |        |   |
            ---          ---          ---          ---

[参照：p.62]

---

4. 小規模な図。

表現すべき命題。

（注：原著の「Smaller Diagram」は「小規模な図」と訳しました。）

```
-----------------
|     |     |
|     x     |
|--y--|--y'--|
|     x'    |
|     |     |
-------------

1. すべての y は x である。
2. 一部の y は x でない。
3. 一切の非-x は非-y ではない。
4. 一部の x は y でない。
5. 一部の非-y は x である。
6. 一切の非-x は y ではない。
7. 一部の非-x は非-y である。
8. 一切の非-x は非-y である。
9. 一部の非-y が存在する。
10. 一切の非-x が存在しない。
11. 一部の y は x であり、また一部は x でない。
12. すべての x は y であり、すべての非-y は非-x である。

[参照：p.62,3]

「国々」を宇宙とし、「文明」とする x、「戦乱」を表す y

13. 一切の未文明国の国々が戦乱的ではない。
14. すべての非戦乱的な国々は未文明である。
15. 一部の国々は非戦乱的である。
16. すべての戦乱的な国々は文明であり、またすべての文明国々は戦乱的である。

17. 一切の国が未文明ではない。

-----------------
「カrocodile」を宇宙とし、「飢餓」とする x、「好ましい」を表す y

18. すべての飢餓的なカrocodileが非好ましいである。
19. 一切のカrocodileが飢餓的であるとき好ましくない。
20. 飢餓でない一部のカrocodileは好ましい；しかし、それらのうち一部はそうではない。
21. 一切のカrocodileが好ましくなく、また一部は飢餓的である。
22. すべてのカrocodileが飢餓的でないとき好ましいであり、またすべての非好ましいカrocodileは飢餓的である。

23. 飢餓的な一部のカrocodileは好ましく、また一部の非飢餓的なカrocodileは非好ましい。

[参照：p.63,4]

5. 小規模な図形

記号の解釈

-----------------
|     |     |
|     x     |
|--y--|--y'--|
|     x'    |
|     |     |
-------------

-----------------
|   |   |           |   |   |
1.  |---|---|       2.  |---|---|
| 1 |   |           |   | 0 |

-----------------
|   | 1 |           |   |   |
3.  |---|---|       4.  |---|---|
|   | 0 |           | 0 | 0 |

```

---

「家」を宇宙とし、x="煉瓦で造られた"、y="二階建ての"として、

\[
\begin{array}{c|cc}
 & 0 & \\
\hline
0 & - & \\
\end{array} \quad 
\begin{array}{c|cc}
 & 0 & \\
\hline
0 & - & - \\
\end{array}

\]

\[
\begin{array}{c|cc}
 & 1 & \\
\hline
0 & 0 & \\
\end{array} \quad 
\begin{array}{c|cc}
 & 1 & \\
\hline
0 & 0 & 1 \\
\end{array}

\]

[参照：p. 65]

「男児」を宇宙とし、x="太った"、y="活発な"として、

\[
\begin{array}{c|cc}
 & 1 & 1 \\
\hline
1 & - & \\
\end{array} \quad 
\begin{array}{c|cc}
 & 1 & 1 \\
\hline
0 & 1 & \\
\end{array}

\]

\[
\begin{array}{c|cc}
 & 0 & 1 \\
\hline
1 & - & 1 \\
\end{array} \quad 
\begin{array}{c|cc}
 & 0 & 1 \\
\hline
0 & 1 & 1 \\
\end{array}

\]

\_\_\_

「猫」を宇宙とし、x="緑色の目を持つ"、y="優れた性質を持つ"として、

\[
\begin{array}{c|cc}
 & 0 & 0 \\
\hline
0 & - & \\
\end{array} \quad 
\begin{array}{c|cc}
 & 0 & 0 \\
\hline
1 & - & \\
\end{array}

\]

\[
\begin{array}{c|cc}
 & 1 & \\
\hline
0 & 1 & \\
\end{array} \quad 
\begin{array}{c|cc}
 & 1 & \\
\hline
0 & 1 & 0 \\
\end{array}

\]

[参照：p. 65, p. 6]

\_\_\_

**より大きな図**

表現すべき命題。

\_\_\_

\[
\begin{array}{c|ccc}
 & \text{--}x\text{--} & | & \text{|  |  |} \\
\hline
\text{| --y--m--y'--} & | & \text{|  |  |} \\
\hline
 & \text{--}x'\text{--} & 
\end{array}
\]

\_\_\_

1. xはmでない。
2. yの一部はm'である。
3. 全てのmはx'である。
4. m'はy'でない。
5. mはxでない；全てのyはmである。
6. xの一部はmであり、yはmでない。

---

7. 全ての m は x' である；m の一切が y でない。

8. x' の一切が m でない；y' の一切が m' でない。

（参照：第67、8ページ）

「うさぎ」を宇宙とし、「m=" greedy "」、「x=" old "」、「y=" black "」として表すと、

9. すべての古いうさぎは greedy ではない。

10. 一部の greedy でないうさぎが黒い。

11. 全ての白いうさぎは greedy でない。

12. 全ての greedy なうさぎは若い。

13. すべての古いうさぎは greedy ではない；すべての黒いうさぎは greedy である。

14. すべての greedy でないうさぎが黒い；すべての古いうさぎは greedy でない。

\_\_\_\_\_\_\_\_\_\_

「鳥」を宇宙とし、「m=" loud singing "」、「x=" well-fed "」、「y=" happy "」として表すと、

15. すべてのwell-fedな鳥はloud singingである；loud singingの一切の鳥が不快ではない。

16. loud singingでない鳥の一切が不快である；すべてのwell-fedな鳥がloud singingである。

\_\_\_\_\_\_\_\_\_\_

「人間」を宇宙とし、「m=" in the house "」、「x=" John "」、「y=" tooth-ache がある "」として表すと、

17. Johnは家にいる；家の中のすべての人間は歯痛で苦しんでいる。

18. 家の中にいるのはJohnだけである；家外の人間は歯痛がない。

\_\_\_\_\_\_\_\_\_\_

（参照：第68-70ページ）

「人間」を宇宙とし、「m=" I "」、「x=" a walk をした "」、「y=" better になった "」として表すと、

19. 私は散歩に出た；私はとても元気になっている。

\_\_\_\_\_\_\_\_\_\_

あなた自身の「宇宙」などを選んで、表すと

20. 私は彼に猫を連れてくるように遣った；彼は kettle を私に持ってきてしまった。

（参照：第70ページ、1）

\_\_\_\_\_\_\_\_\_\_

注意：各問題に対して、x と y のみを対象とする小さな図を描き、与えられた大きな図に従ってマークし、その後、この小さな図からできる限りの命題を読み取るべきである。

---

-----------              -----------
              |0    |     |            |     |     |
              |   --|--   |            |   --|--   |
              |  |0 | 0|  |            |  |0 | 1|  |
          1.  |--|--|--|--|        2.  |--|--|--|--|
              |  |1 |  |  |            |  |0 |  |  |
              |   --|--   |            |   --|--   |
              |0    |     |            |     |     |
               -----------              -----------

[参照：p.72]

-----------              -----------
              |     |     |            |     |    0|
              |   --|--   |            |   --|--   |
              |  |0 | 0|  |            |  |  |  |  |
          3.  |--|--|--|--|        4.  |--|--|--|--|
              |  |1 | 0|  |            |  |0 |  |  |
              |   --|--   |            |   --|--   |
              |     |     |            |     |    0|
               -----------              -----------

\_\_\_\_\_\_\_\_\_\_

大図形において、前節の以下の命題のペアをマークし、それに従って小図形を描き、そのような操作を繰り返す。

5. No.13 [参照：p.49]    9. No.17
6. No.14                 10. No.18
7. No.15                 11. No.19 [参照：p.50]
8. No.16                 12. No.20

\_\_\_\_\_\_\_\_\_\_

大図形において、以下の命題のペアをマークし、それに従って小図形を描き、そのような操作を繰り返す。これらは実際には三段論法の前提のペアであり、小図形から読み取った結果が結論である。

13. 兴奮する本は熱い患者に適さない；興奮しない本は人を眠らせる。
14. 些か、公正な報酬を得る者はその公正な報酬を得ている；勇敢以外の者は公正な報酬を得ない。
15. 子供は忍耐強くない；忍耐強くない者は静かに座れない。

[参照：p.72-5]

16. すべての豚は太い；すべての骨格は太くない。
17. すべての猿は兵士ではない；すべての猿は悪戯好きである。
18. 私の諸兄妹のうち誰も正直者ではない；すべての判事は不正直者である。
19. 些か、雨の日がある；雨の日は退屈である。
20. すべての薬は苦い；センナは薬である。
21. 些か、ユダヤ人は裕福である；すべてのパタゴニア人はゲンティルである。
22. すべての禁酒者は砂糖を好む；夜莺はワインを飲まない。

---

23. もののパンケーキは健康に良いものではない；すべてのバウムクーヘンは健康に良くないものである。

24. ものの太った動物は走れない；一部のグレイハoundsは走れる。

25. すべての兵士は歩行する；一部の若者は兵士でない。

26. 砂糖は甘い；塩は甘くない。

27. 一部の卵は硬煮えである；すべての卵は破けにくいものではない。

28. 家にはユダヤ人はいない；庭にはゲーテンリック人はいない。

（参照：p. 75-82）

29. すべての戦いは騒々しい；騒がしくないものは注意を引かない。

30. ゲーテンリック人には狂人はいない；すべてのラビはゲーテンリック人である。

31. すべての魚は泳げる；一部のスケートは魚である。

32. すべての情熱的な人々は理不尽である；一部の演説家は情熱的である。

（参照：p. 82-84）

\newpage

\chapter{歪んだ答え（Crooked Answers）}
\label{ch:iii}

\begin{center}
\large\textit{三段論法とダイアグラム配置の詳細解答と記号論理的解法}
\end{center}
\vspace{4mm}

第三章

曲がった答え

「私は彼に、自分の思うように答えた。
“紅いさばが木からどれだけ生えるかのように、その数だ。”」

\_\_\_\_\_\_\_\_\_\_\_\_

1. 基礎的な定義。

1. 何かを「ある物に属する」と言えるとき、つまり「ある物に属するもの」と言えるとき、それを「属性」と呼ぶ。例えば、「焼いた」は「パン」に頻繁に属し、「美しい」は「赤ちゃん」に希少に属する。

2. 二つの物の名前（例：「これらの豚は太い動物である」）や二つの属性の名前（例：「ピンクは明るい赤色である」）で表されるとき。

3. 物の名前と属性の名前の組み合わせ（例：「これらの豚はピンクである」）。物は実際には属性を有することはできない。

4. 意図的に主語の名詞が文末に繰り返されること（例：「これらの豚はピンクである（豚）」）。

5. 「命題」は、あるクラスのすべてまたは一部の物が別のクラスの物であることを述べる文である。このクラスを「主項」と呼び、その他のクラスを「述項」と呼ぶ。例えば、「新しいケーキのいくつかは美味しいわけではない」これは完全に書くと「新しいケーキのいくつかは美味しいではないケーキである」となり、クラス「新しいケーキ」が主項であり、クラス「美味しいでないケーキ」が述項である。

6. 主項の一部の物が特定の状態にあることを述べる命題を「個別命題」と呼ぶ。例えば、「新しいケーキのいくつかは美味しい」「新しいケーキのいくつかは美味しいではない」。

7. 主項のすべてまたは一部の物が特定の状態にあることを述べる命題を「普遍命題」と呼ぶ。例えば、「新しいケーキのすべては美味しいわけではない」「新しいケーキのすべては美味しいではない」。

8. 各 compartment 内の物には、その compartment の二つの辺に記載されている符号で表される二つの属性が存在する。

9. 全体的な図形が属するクラスの名前として使用される。

10. 二つの声明を含む命題。例えば、「新しいケーキのいくつかは美味しいし、いくつかは美味しくない」。

11. 全体が、その全体を分割した集合群に「尽きる」場合、つまり、どの要素もこれらのいずれかの集合に属する場合がある。例えば、「新ケーキ」のクラスは、「美味しい」と「美味しいでない」に「完全に」分割されているため、すべての新しいケーキはどちらか一方である。

12. ある人が二つの派閥のどちらに加入すべきか決断できず、その状態を「柵に座っている」と表現し、どの一方にも決定できない状況を示す。

13. 「一部 x は y」および「すべての x は y'」。

14. 主語が一つの事物である命題を「個別命題」と呼ぶ。例えば、「私は幸せです」「ジョンは家にいない」は個別命題であり、これらは「存在する全ての私たちは幸せである」と同じ意味を持つ普遍的命題と等価である。

15. 「一部」または「すべて」で始まる命題。

16. これらの命題が「一部」または「ない」で始まることを指す。例えば、「一部 abc は def」は「一部 bf は acde」と再配置され、これらは「一部 abcdef が存在する」と等価である。

17. 一部のトラウトは凶暴であり、すべてのトラウトは凶暴でない。

18. 一部の茹でた卵は不健康であり、すべての茹でた卵は健康でない。

19. 一部の I が幸せであり、すべての I は不幸せではない。

20. 一部のジョンは家にいないであり、すべてのジョンは家にいるではない。

21. 大きな図形の各 compartiment 内の事物には、THREE つの属性があり、これらの符号は compartiment の三隅（m' の場合を除く）に記載されている。

22. 物事が三つの異なる属性について分割され、二つの命題が与えられ、これら二つの命題は異なる属性の組み合わせを含むとし、さらに、これらの命題からまだ同時に現れていなかった二つの属性を含む第三の命題が証明できる場合、その命題は「前提」であり、第三の命題は「結論」とされ、これら全てが「三段論法」と呼ばれる。例えば、「すべて m は x' でない」と「すべて m' は y」が前提となり、それらから x と y を含む結論を証明できる場合がある。

23. もし前提に共通して出現する属性がある場合、その属性を含む語句は「中間語句」と呼ばれる。例えば、「一部のmはxである」および「すべてのmはyではない」が前提の場合、「m事物」のクラスは「中間語句」とされる。

もし一つの前提に属性が出現し、もう一つの前提にはその属性の否定形が現れている場合、それらを含む語句は「中間語句」と呼ぶことができる。例えば、「すべてのmはxではない」および「すべてのm'はyである」が前提の場合、「m事物」のクラスと「m'事物」のクラスは「中間語句」と呼ばれる。

24. これらは確実に識別できるため、肯定命題（すなわち、「一部」または「すべて」で始まるもの）は時として赤い駒を「境界線上に置く必要がある」という状況を要求することがある。

25. 我々が心配するのは結論が前提から論理的に導かれるかどうかである。したがって、もし前提が真であれば、結論も真であることが確認される。

26. 赤い駒は「この compartiment は占有可能である」という意味であり、灰色の駒は「この compartiment は占有不可能である」または「この compartiment は必ず空である」という意味であると理解する。

27. 「論理的な誤りのある前提」および「論理的な誤りのある結論」。

28. 大きな図形から小さな図形へのマークの移動を試みた際に、その四つの compartiment のいずれにも情報がない場合。

29. 正しい結論を見つけてから、提供された結論が正確な結論と同一であるか部分的に含まれているか確認する。

30. 提供された結論が正しい結論の一部の場合。この場合、その結論は「不完全な結論」と呼ばれることがある。

2.
小さな図形の半分を表す。

命題：

\_\_\_\_\_\_\_\_\_\_

|   |   |
1.  |   | 1 |
|   |   |

|   |   |
2.  | 0 | 1 |
|   |   |

|   |   |
3.  | 1 | 1 |
|   |   |

|   |   |
4.  | 0 | 0 |
|   |   |

-----------------
|   |   |        |   |   |
5.  |   1   |      6.  |   | 0 |
|   |   |        |   |   |
-----------------        -----------------

-----------------
|   |   |
7.  | 1 | 1 | ある x が存在することを表現するために
|   |   |
-----------------        -----------------
          |   |   |
       |   1 1 | のような図形が考えられるかもしれません。
          |   |   |
           -------
Diagramは「ある x が y である」という文に含まれている実質的な内容を表しています。 compartment の一つが占拠されていることを教えてくれる赤い駒を置くことは、すでに ONE が占拠されていることを知っている限り、既に分かっていることになります。
           -----------------
          |   |   |
8.  No x are y. すなわち | 0 |   |
|   |   |
-----------------

-----------------
           |   |   |
9.  Some x are y'. すなわち |   | 1 |
|   |   |
----------------

-----------------
           |   |   |
10. All x are y. すなわち | 1 | 0 |
|   |   |
----------------

-----------------
           |   |   |
11. Some x are y. すなわち | 1 |   |
|   |   |
----------------

-----------------
           |   |   |
12. No x are y. すなわち | 0 |   |
|   |   |
----------------

-----------------
                |   |   |
13. Some x are y, and some are y'. すなわち | 1 | 1 |
|   |   |
-----------------

-----------------
           |   |   |
14. All x are y'. すなわち | 0 | 1 |
|   |   |
----------------

---
             |   |
15. No y are x'. すなわち |---|
| 0 |
-------------

\#\#\# 16. 全ての y は x である。すなわち、
\[
\begin{array}{c}
\hline
1 \\
\hline
0 \\
\end{array}
\]

\#\#\# 17. y の存在しないものである。すなわち、
\[
\begin{array}{c}
\hline
0 \\
\hline
0 \\
\end{array}
\]

\#\#\# 18. 一部の y は x' である。すなわち、
\[
\begin{array}{c}
\hline
\phantom{1} \\
\hline
1 \\
\end{array}
\]

\#\#\# 15. 一部の y が存在する。すなわち、
\[
\begin{array}{c}
\hline
-1 \\
\hline
\phantom{1} \\
\end{array}
\]

---

\#\#\# 3. 小図形の半分。

記号の解釈。

\_\_\_\_\_\_\_\_\_\_

1. 全ての x は y' でない。
2. x の存在しないものである。
3. 一部の x が存在する。
4. 全ての x は y' である。
5. 一部の x は y である。すなわち、一部の良い謎は難しくなる。
6. 全ての x は y である。すなわち、全ての良い謎は難しくなる。
7. x の存在しないものである。すなわち、riddles の存在しないものである。
8. 全ての x は y でない。すなわち、全ての良い riddles は難しくない。
9. 一部の x は y' である。すなわち、一部のクラブは無自発的である。
10. 全ての x は y でない。すなわち、全てのクラブは自己中心的なものである。
11. 全ての x は y' である。すなわち、全てのクラブは無自発的である。
12. 一部の x は y であり、一部は y' である。すなわち、一部のクラブは自己中心的なものであり、一部は無自発的である。
13. 全ての y' は x' である。すなわち、全ての不健康者は不幸である。
14. 一部の y' が存在する。すなわち、一部の人々は不健康である。
15. 一部の y' は x であり、一部は x' である。すなわち、一部の不健康者は幸せであり、一部は不幸である。
16. 不健康者の存在しないものである。すなわち、誰もが不健康ではない。

---

\#\#\# 4. 小図形。

命題の表現。

\_\_\_\_\_\_\_\_\_\_

\[
\begin{array}{c|c}
\hline
1 & \\
\hline
0 \\
\end{array}
\quad \text{と} \quad
\begin{array}{c|c}
\hline
0 & \\
\hline
0 \\
\end{array}
\]

1. 
\[
\begin{array}{c|c}
\hline
1 & 0 \\
\hline
0 \\
\end{array}
\]
2.
\[
\begin{array}{c|c}
\hline
0 & 1 \\
\hline
0 \\
\end{array}
\]

3. 
\[
\begin{array}{c|c}
\hline
1 & \phantom{1} \\
\hline
0 \\
\end{array}
\]
4.
\[
\begin{array}{c|c}
\hline
0 & 1 \\
\hline
0 \\
\end{array}
\]

5. 
\[
\begin{array}{|c|c|}
\hline
 & 1 \\
\hline
- & --- \\
 & \text{} \\
\hline
\end{array}
\]

6.
\[
\begin{array}{|c|c|}
\hline
 & 1 \\
\hline
0 & --- \\
 & \text{} \\
\hline
\end{array}
\]

7. 
\[
\begin{array}{|c|c|}
\hline
 & \text{} \\
\hline
- & --- \\
 & 1 \\
\hline
\end{array}
\]

8.
\[
\begin{array}{|c|c|}
\hline
 & \text{} \\
\hline
0 & --- \\
 & 1 \\
\hline
\end{array}
\]

9. 
\[
\begin{array}{|c|c|}
\hline
 & \text{} \\
\hline
- & -1- \\
 & \text{} \\
\hline
\end{array}
\]

10.
\[
\begin{array}{|c|c|}
\hline
 & \text{} \\
\hline
- & --- \\
 & 0 \\
\hline
\end{array}
\]

11.
\[
\begin{array}{|c|c|}
\hline
1 & \text{} \\
\hline
- & --- \\
 & 0 \\
\hline
\end{array}
\]

12.
\[
\begin{array}{|c|c|}
\hline
1 & \text{} \\
\hline
- & --- \\
 & 1 \\
\hline
\end{array}
\]

13. 
\[
\begin{array}{|c|c|}
\hline
 & \text{} \\
\hline
0 & --- \\
 & \text{} \\
\hline
\end{array}
\]
すなわち、「x'は存在しない」。

14.
\[
\begin{array}{|c|c|}
\hline
 & 0 \\
\hline
0 & --- \\
 & 1 \\
\hline
\end{array}
\]
すなわち、「すべてのy'はx'である」。

15. 
\[
\begin{array}{|c|c|}
\hline
 & \text{} \\
\hline
- & -1- \\
 & \text{} \\
\hline
\end{array}
\]
すなわち、「あるy'が存在する」。

16.
\[
\begin{array}{|c|c|}
\hline
1 & 0 \\
\hline
- & --- \\
 & 0 \\
\hline
\end{array}
\]
すなわち、「すべてのyはxであり、すべてのxはyである」。

17.
\[
\begin{array}{|c|c|}
\hline
 & \text{} \\
\hline
0 & --- \\
 & 0 \\
\hline
\end{array}
\]
すなわち、「x'が存在しない」。

18.
\[
\begin{array}{|c|c|}
\hline
0 & 1 \\
\hline
- & --- \\
 & \text{} \\
\hline
\end{array}
\]
すなわち、「すべてのxはy'である」。

---

19. 任意のxがyでない。すなわち、

\[
\begin{array}{|c|c|}
\hline
0 &   \\
\hline
--- & --- \\
\hline
\end{array}
\]

---

20. 一部のx'がyであり、また一部がy'である。すなわち、

\[
\begin{array}{|c|c|}
\hline
   &   \\
\hline
--- & --- \\
\hline
1 & 1 \\
\hline
\end{array}
\]

---

21. yは存在しないし、一部のxが存在する。すなわち、

\[
\begin{array}{|c|c|}
\hline
0 &   \\
\hline
--- & --- \\
\hline
0 &   \\
\hline
\end{array}
\]

---

22. すべてのx'がyであり、またすべてのy'がxである。すなわち、

\[
\begin{array}{|c|c|}
\hline
1 & 0 \\
\hline
--- & --- \\
\hline
1 &   \\
\hline
\end{array}
\]

---

17. 一部のxがyであり、また一部のx'がy'である。すなわち、

\[
\begin{array}{|c|c|}
\hline
   & 1 \\
\hline
--- & --- \\
\hline
1 &   \\
\hline
\end{array}
\]

---

5. 小さい図形。

シンボルの解釈。

\_\_\_\_\_\_\_\_\_\_

1. 一部のyが-xである、または、一部の-xがyである。

2. 任意の-xが-not-yでない、または、任意の-not-yが-not-xでない。

3. 任意の-not-yがxでない。

4. 任意の-not-xは存在しない。すなわち、任意の事物は-not-xではない。

5. yは存在しない。すなわち、二階建ての家は存在しない。

6. 一部のx'が存在する。すなわち、一部の家は煉瓦で作られていない。

7. xはy'でない。または、y'はxでない。すなわち、煉瓦で作られた家は二階建て以外のものではない。または、二階建てでない家は煉瓦で作られていない。

8. すべてのx'がy'である。すなわち、煉瓦で作られていない家は二階建てでない。

9. 一部のxがyであり、また一部がy'である。すなわち、一部の太い男の子は活発であり、また一部はそうではない。

10. すべてのy'がx'である。すなわち、すべての怠け者の男の子は痩せている。

11. すべてのxがy'であり、またすべてのy'がxである。すなわち、すべての太い男の子は怠け者であり、すべての怠け者は太い。

12. すべてのyがxであり、またすべてのx'がyである。すなわち、すべての活発な男の子は太く、すべての痩せたものは怠け者である。

---

13. xは存在しないし、y'も存在しない。すなわち、猫には緑色の目がないし、悪い性質を持つ猫もない。

14. あるxがy'であり、あるx'がyである。または、あるyがx'であり、あるy'がxである。すなわち、緑色の目を持った猫は悪い性質を持っているし、それ以外の猫の中にも良い性質を持つ猫がある。または、良い性質を持つ猫の中に緑色の目がない猫もいるし、悪い性質を持つ猫の中に緑色の目を持つ猫もいる。

15. あるxがyであり、あるx'がy'でない。または、あるyがxであり、あるy'がx'でない。すなわち、緑色の目を持った猫の中には良い性質を持つ猫があるし、それ以外の猫の中には悪い性質を持つ猫がない。または、良い性質を持つ猫の中に緑色の目を持つ猫もいるし、悪い性質を持つ猫の中に緑色の目を持っていない猫がいない。

16. すべてのxがy'であり、すべてのx'がyである。または、すべてのyがx'であり、すべてのy'がxである。すなわち、すべての緑色の目を持った猫は悪い性質を持っているし、それ以外の猫は全て良い性質を持っている。または、全ての良い性質を持つ猫には緑色の目がないし、全ての悪い性質を持つ猫には緑色の目がある。

6. 大きな図面。

命題を表す。

\_\_\_\_\_\_\_\_\_\_

---------------               ---------------
            |       |       |             |       |       |
            |    ---|---    |             |    ---|---    |
            |   | 0 | 0 |   |             |   |   |   |   |
        1.  |---|---|---|---|         2.  |-1-|---|---|---|
            |   |   |   |   |             |   |   |   |   |
            |    ---|---    |             |    ---|---    |
            |       |       |             |       |       |
             ---------------               ---------------

---------------               ---------------
            |       |       |             |       |     0 |
            |    ---|---    |             |    ---|---    |
            |   | 0 | 0 |   |             |   |   |   |   |
        3.  |---|---|---|---|         4.  |---|---|---|---|
            |   |   -   |   |             |   |   |   |   |
            |    ---|---    |             |    ---|---    |
            |       |       |             |       |     0 |
             ---------------               ---------------

以上が原著の

5. 
\[
\begin{array}{|c|c|c|}
\hline
0 & 0 & 0 \\
\hline
0 & 1 & 1 \\
\hline
\end{array}
\]

6. 
\[
\begin{array}{|c|c|c|c|}
\hline
0 & 0 & 0 & 0 \\
\hline
0 & 1 & 1 & 1 \\
\hline
0 & 0 & 0 & 0 \\
\hline
\end{array}
\]

7. 
\[
\begin{array}{|c|c|c|}
\hline
& & 0 \\
\hline
--- & --- & --- \\
\hline
0 & 0 & 0 \\
\hline
\end{array}
\]

8. 
\[
\begin{array}{|c|c|c|c|}
\hline
& & & 0 \\
\hline
--- & --- & --- & --- \\
\hline
0 & 0 & 1 & 0 \\
\hline
--- & --- & --- & --- \\
\hline
& & & 0 \\
\hline
\end{array}
\]

9. 
\[
\begin{array}{|c|c|c|}
\hline
& & \\
\hline
--- & --- & \\
\hline
0 & 0 & 0 \\
\hline
\end{array}
\]
「ある x は m ではない。すなわち、」
\[
\begin{array}{|c|c|c|c|}
\hline
--- & --- & --- & --- \\
\hline
0 & 0 & 0 & 0 \\
\hline
0 & 1 & 1 & 1 \\
\hline
\end{array}
\]

10. 
\[
\begin{array}{|c|c|c|}
\hline
& & \\
\hline
--- & --- & \\
\hline
0 & 0 & 0 \\
\hline
\end{array}
\]
「ある m' は y である。すなわち、」
\[
\begin{array}{|c|c|c|c|}
\hline
-1 & --- & --- & --- \\
\hline
0 & 0 & 0 & 0 \\
\hline
--- & --- & --- & \\
\hline
& & & \\
\hline
\end{array}
\]

11. 全ての y' は m' である。すなわち、

\[
\begin{array}{|c|c|c|c|}
\hline
 &  & 0 & \\
\hline
-1 & - & 0 & - \\
\hline
\end{array}
\]

12. 全ての m は x' である。すなわち、

\[
\begin{array}{|c|c|c|c|}
\hline
 &  & 0 & 0 \\
\hline
- & - & - & 1 \\
\hline
\end{array}
\]

13. 全ての x は m でない；すなわち、全ての y は m である。

\[
\begin{array}{|c|c|c|c|}
\hline
0 &  &  & \\
\hline
- & - & 1 & - \\
\hline
 &  &  & 0 \\
\hline
\end{array}
\]

14. 全ての m' は y である；すなわち、全ての x は m でない。

\[
\begin{array}{|c|c|c|c|}
\hline
0 &  & 0 & \\
\hline
- & - & - & - \\
\hline
1 &  & 0 & 0 \\
\hline
\end{array}
\]

15. 全ての x は m である；すなわち、全ての m は y' である。

\[
\begin{array}{|c|c|c|c|}
\hline
0 &  & 0 & \\
\hline
- & - & 1 & 0 \\
\hline
 &  & 0 & 0 \\
\hline
\end{array}
\]

16. 全ての m' は y'；すなわち、全ての x が m' でない。
   即ち、No x が m'。

17. 全ての x が m；すなわち、全ての m が y。
   [参照：第7番注記、第60ページ。]

18. 全ての x' が m でない；すなわち、全ての m' が y'。
   
19. 全ての m が x；すなわち、全ての m が y。

20. 我々は「人」を宇宙と定義すべきである。私自身を「中項」とする場合、前提は以下のようになる：

私は誰かに猫を持ってくるように送った人間である；
   私は誰かが誤って kettle を持ってきた人間に会った人間である。

または、「彼」を中項とする場合、前提は以下のようになる：

彼は私が猫を持ってきて欲しかった人間である；
   彼は私が誤って kettle を持って来た人間である。

後者の形式が最も適切であると考えられる。この短編の興味は、彼の愚かさに依存しているからであり、私自身の経験とは無関係である。したがって、

m = 「彼」; x = 「私が送った人々など」; y = 「持ってきた人々など」

となる。従って、以下の図表が必要である。

\[
\begin{array}{|c|c|}
\hline
 &  \\
\hline
 & \text{---|---} \\
\hline
 & |1|0|   \\
\hline
\text{---|---|---|---} \\
\hline
 & |0|0|   \\
\hline
 & \text{---|---} \\
\hline
 &  \\
\hline
\end{array}
\]

7. 使用する図表は両方である。

\[
   \begin{array}{|c|c|}
   \hline
   0 |   \\
   \hline
   \text{|---|---|} & \text{i.e., 全ての y は x' である。} \\
   1 |   \\
   \hline
   \end{array}
   \]

\[
   \begin{array}{|c|c|}
   \hline
   |   | 1 \\
   \hline
   \text{|---|---|} & \text{i.e., 一部の x は y' である。または、一部の y' は x である。} \\
   |   |   \\
   \hline
   \end{array}
   \]

\[
   \begin{array}{|c|c|}
   \hline
   |   |   \\
   \hline
   \text{|---|---|} & \text{i.e., 一部の y は x' である。または、一部の x' は y である。} \\
   | 1 |   \\
   \hline
   \end{array}
   \]

\[
   \begin{array}{|c|c|}
   \hline
   |   |   \\
   \hline
   \text{|---|---|} & \text{i.e., 一部の x' は y' である。または、一部の y' は x である。} \\
   |   | 0 \\
   \hline
   \end{array}
   \]

\[
   \begin{array}{|c|c|}
   \hline
   0 |   \\
   \hline
   \text{|---|---|} & \text{i.e., 全ての y は x' である。つまり、全ての黒い兎は若々しい。} \\
   1 |   \\
   \hline
   \end{array}
   \]

\[
   \begin{array}{|c|c|}
   \hline
   |   |   \\
   \hline
   \text{|---|---|} & \text{i.e., 一部の y は x' である。つまり、一部の黒い兎は若々しい。} \\
   | 1 |   \\
   \hline
   \end{array}
   \]

\[
   \begin{array}{|c|c|}
   \hline
   1 | 0 \\
   \hline
   \text{|---|---|} & \text{i.e., 全ての x は y である。つまり、全ての肥満した鳥は幸せである。} \\
   |   |   \\
   \hline
   \end{array}
   \]

\[
   \begin{array}{|c|c|}
   \hline
   |   |   | & \text{i.e., 一部の x' は y' である。つまり、一部の未肥満した鳥は不幸である。} \\
   \hline
   |---|---| & \text{または、一部の不幸な鳥は未肥満である。} \\
   |   | 1 | \\
   \hline
   \end{array}
   \]

\[
   \begin{array}{|c|c|}
   \hline
   1 | 0 \\
   \hline
   \text{|---|---|} & \text{i.e., 全ての x は y である。つまり、ジョンは歯痛を持っている。} \\
   |   |   \\
   \hline
   \end{array}
   \]

\[
   \begin{array}{|c|c|}
   \hline
   |   |   | & \text{i.e., 一部の x' は y である。つまり、一部の未肥満した鳥は幸せである。} \\
   \hline
   |---|---| & \text{または、一部の幸せな鳥は未肥満である。} \\
   |   | 1 | \\
   \hline
   \end{array}
   \]

\[
   \begin{array}{|c|c|}
   \hline
   1 | 0 \\
   \hline
   \text{|---|---|} & \text{i.e., 全ての x は y である。つまり、ジョンだけが歯痛を持っている。} \\
   |   |   \\
   \hline
   \end{array}
   \]

---
| 1 |   |
11. |---|---| 即ち、ある x は y である。即ち、ある人が散歩をした後で気分が良くなった。
   -------

---
| 1 |   |
12. |---|---| 即ち、ある x は y である。即ち、私が猫を連れてくるように送った人間が、鍋を間違えて持ってきてしまった。
   -------

---
|       |     0 |
|    ---|---    |
|   | 0 | 0 |   
13. |-1-|---|---|---|
   |   |   |   |
   |    ---|---    |
   |       |     0 |
   ---------------
      即ち、以下のように定義します：
         "本"を宇宙とし、m="興奮する"、x="熱のある患者に適したもの"、y="人を眠らせてしまうもの"とする。
         
         m のすべては x でない；従って y' のすべても x でない。
         m' のすべては y である。

即ち、熱のある患者に適さない本は、人を眠らせることのできるもの以外はない。

---------------
|       |       |
|    ---|---    |
|   | 1 | 0 |   
14. |---|---|---|---|
   |   |   | 0 | 
   |    ---|---    |
   |       |       |
   ---------------
      即ち、以下のように定義します：
         "人間"を宇宙とし、m="公正に扱われるべき人"、x="報酬を受け取る人"、y="勇敢な人"とする。
         
         m の一部は x である；従って y の一部も x である。
         y' のすべては m でない。

即ち、勇敢な人々の一部が報酬を受け取る。

---------------
| 0     |       |
|    ---|---    |
|   | 0 | 0 |   
15. |---|---|---|---|
   |   |   |   |
   |    ---|---    |
   | 0     |       |
   ---------------

即ち、以下のように定義します：
         "人間"を宇宙とし、m="患者"、x="子供"、y="静かに座れるもの"とする。
         
         x のすべては m でない；従って x のすべても y でない。
         m' のすべても y でない。

即ち、子供のすべてが静かに座ることはできない。

---------------
| 0     |     0 |
|    ---|---    |
|   | 0 | 1 |   
16. |---|---|---|---|
   |   | 0 |   |
   |    ---|---    |
   |       |       |
   ---------------

即ち、以下のように定義します：
         "人間"を宇宙とし、m="患者"、x="子供"、y="静かに座れるもの"とする。
         
         x のすべては m でない；従って x のすべても y でない。
         m' のすべても y でない。

即ち、子供のすべてが静かに座ることはできない。

「事物」を宇宙とし、m="太い"; x="豚"; y="骨格"とする。

すべてのxはmである；従って、すべてのxはy'である。
yはmでない。

すなわち、すべての豚は骨格ではない。

-------------------
|       |       |
|    ---|---    |
|   | 0 | 0 |   |
17.  |---|---|---|---|
|   | 1 | 0 |   |
|    ---|---    |
|       |       |   |
-------------------

「生物」を宇宙とし、m="猿"; x="兵士"; y="悪戯好き"とする。

mはxでない；従って、一部のyはx'である。
すべてのmはyである。

すなわち、一部の悪戯好きな生物は兵士ではない。

-------------------
| 0     |       |
|    ---|---    |
|   | 0 | 0 |   |
18.  |---|---|---|---|
|   |   |   |   |
|    ---|---    |
| 0     |       |
-------------------

「人間」を宇宙とし、m="正直な"; x="私の従兄弟"; y="裁判官"とする。

xはmでない；従って、xはyでもない。
yはm'である。

すなわち、私の従兄弟の誰もが裁判官ではない。

-------------------
|       |       |
|    ---|---    |
|   | 1 | 0 |   |
19.  |---|---|---|---|
|   |   |   |   |
|    ---|---    |
|       |       |
-------------------

「期間」を宇宙とし、m="日"; x="雨の"; y="疲れる"とする。

一部のmはxである；従って、一部のxはyである。
すべてのxmはyである。

すなわち、一部の雨の期間は疲れる。

注記：これらは正当な前提ではない。結論は実際には第二前提の一部であり、第一前提は余分なものだからである。これは文字で示すと次のようになる：

「すべてのxmはyである」に含まれる「一部のxmはyである」が、「一部のxはyである」という結論を含んでいるからである。または、言葉で言えば、「すべての雨の日は疲れる」に含まれる「一部の雨の日は疲れる」が、「一部の雨の期間は疲れる」という結論を含んでいるからである。

さらに、第一前提は余分であり、実際には第二前提に含まれている。なぜならそれは「雨の日が存在する」ということと等価であり、これは我々が知るところでは「すべての雨の日は退屈である」から導かれる事柄だからである。

全体として、非常に不十分な一組の前提である！

---------------
     | 0     |       |
     |    ---|---    |
     |   | 1 |   |   |
20.  |---|---|---|---|        -------
     |   | 0 | 0 |   |       | 1 |   |
     |    ---|---    |       |---|---|
     | 0     |       |       | 0 |   |
      ---------------

「事物」を宇宙とし、m="薬"; x="不快な"; y="セナ"とする。

すべてのmはxである；従って、すべてのyはxである。
すべてのyはmである。

すなわち、「セナは不快なものである」ということである。

[参照：第7番注、60ページ]

---------------
     |       |       |
     |    ---|---    |
     |   | 0 | 1 |   |
21.  |-1-|---|---|---|        -------
     |   | 0 |   |   |       |   | 1 |
     |    ---|---    |       |---|---|
     |       |       |       |   |   |
      ---------------

「人」を宇宙とし、m="ユダヤ人"; x="裕福な"; y="パタゴニア人"とする。

一部のmはxである；従って、一部のxはy'である。
すべてのyはm'である。

すなわち、「一部の裕福な人はパタゴニア人ではない」ということである。

---------------
     | 0     |       |
     |    ---|---    |
     |   |   -   |   |
22.  |---|---|---|---|        -------
     |   | 0 | 0 |   |       |   |   |
     |    ---|---    |       |---|---|
     | 0     |       |       | 0 |   |
      ---------------

「生物」を宇宙とし、m="禁酒者"; x="砂糖が好む"; y="夜莺"とする。

すべてのmはxである；従って、yの一部はx'である。
yの一部はm'である。

すなわち、「夜莺の一部は砂糖を好まない」ということである。

---------------
     |       |       |
     |    ---|---    |
     |   | 0 | 0 |   |
23.  |-1-|---|---|---|        -------
     |   | 0 |   |   |       |   |   |
     |    ---|---    |       |---|---|
     |       |       |       |   |   |
      ---------------

「食事」を宇宙とし、m="栄養価のある"; x="マフィン"; y="パン"とする。

---

あるxはmでない；
すべてのyはmである。

この小さな図には「情報がない」ので、結論を導くことはできない。

-------------------
|       |       |
|    ---|---    |
|   | 0 | 0 |   |
24.  |---|---|---|---|
|   | 1 |   |   |
|    ---|---    |
|       |       |
-------------------
|       |       |
|    ---|---    |

令「生物」を宇宙とし、「m＝走りが良い」とする；「x＝太い」、「y＝グレイハounds」である。

したがって、あるxはmでない；従って、あるyはx'である。
あるyはmである。

すなわち、あるグレイハoundsは太くない。

-------------------
|       |       |
|    ---|---    |
|   |   -   |   |
25.  |-1-|---|---|---|
|   | 0 | 0 |   |
|    ---|---    |
|       |       |
-------------------
|       |       |
|    ---|---    |

令「人間」を宇宙とし、「m＝兵士」とする；「x＝歩行者」、「y＝若者」である。

すべてのmはxである；
あるyはm'である。

この小さな図には「情報がない」ので、結論を導くことはできない。

-------------------
| 0     |     0 |
|    ---|---    |
|   | 0 | 1 |   |
26.  |---|---|---|---|
|   | 0 |   |   |
|    ---|---    |
| 1     |       |
-------------------
|       |       |
|    ---|---    |

令「食料」を宇宙とし、「m＝甘い」とする；「x＝砂糖」、「y＝塩」である。

すべてのxはmである；従って、すべてのxはy'である。
すべてのyはm'である。従って、すべてのyはx'である。

すなわち、砂糖は塩でない。
塩は砂糖でもない。

-------------------
|       |       |
|    ---|---    |
|   | 1 | 0 |   |
27.  |---|---|---|---|
|   |   | 0 |   |
|    ---|---    |
|       |       |
-------------------
|       |       |
|    ---|---    |

令「物」を宇宙とし、「m＝卵」とする；「x＝茹でたもの」、「y＝割れるもの」である。

あるmはxである；従って、あるxはyである。
あるmはy'でない。

すなわち、ある茹でたものは割れる可能性がある。

28. 人を全宇宙とする。m="ユダヤ人"、x="家にいるもの"、y="庭にいるもの"とする。

mはxでない；従って、xはyでない。
             m'はyでない。

すなわち、家にいる人たちは庭にもいない。

---------------
     | 0     |       |
     |    ---|---    |
     |   | 0 | 0 |   |
29. 事物を全宇宙とする。m="騒々しい"、x="戦闘"、y="見落としができるもの"とする。

xはすべてmである；従って、x'の一部はyである。
             m'はすべてyである。

すなわち、戦闘でない事物の一部は見落としができる。

---------------
     | 0     |       |
     |    ---|---    |
     |   | 0 | 0 |   |
30. 人を全宇宙とする。m="狂人"、x="拉比"とする。

mはxでない；従って、すべてのyはx'である。
             すべてのyはmである。

すなわち、すべての拉比は sane（冷静な）である。

---------------
     |       |       |
     |    ---|---    |
     |   | 1 |   |   |
31. 物事を全宇宙とする。m="魚"、x="泳げる"、y="スケート"とする。

mはすべてx'である；従って、yの一部はxである。
             yの一部はmである。

すなわち、スケートの一部は泳げることがある。

32. 宇宙を「人間」、mを「情熱的な人」、xを「合理的な人」、yを「演説家」と定義する。

すべてのmはx'であり；従って、一部のyはx'である。
一部のyはmである。

すなわち、「一部の演説家は合理的でない」ということである。

なお、7番目の問題に関する注を参照されたい。（第60頁）

\newpage

\chapter{当たるも八卦当たらぬも八卦（Hit or Miss）}
\label{ch:iv}

\begin{center}
\large\textit{総合論理パズル・誤謬の看破とゲームとしての記号論理}
\end{center}
\vspace{4mm}

第四章

命中または失敗

「汝はそれを打て、打て、打て。
汝はそれを打つことができない、よしんばの男よ。」

\_\_\_\_\_\_\_\_\_\_

1. 痛みは疲れるものである；痛みのないものは誰もが喜んで望まない。

2. 金髪の人は髪ブラシを必要としない；爬虫類には髪はない。

3. 考えのない人々は悪事を成す；考えの深い人は約束を忘れない。

4. ジョンを好まない；私の友人の中にはジョンを好む者がいる。

5. じゃがいもはパイナップルではない；すべてのパイナップルは美味しい。

6. 鋤は野心的でない；針は鋲でもない。

7. 私のすべての友人は風邪を持っている；風邪を患っている者は誰も歌うことができない。

8. これらのすべての料理はよく調理されている；調理が不十分な場合にのみ不健康な料理がある。

9. 薬は誰もが好まない；セナは薬である。

10. 一部のオイスターは静かである；静かな生物は誰もが面白くない。

11. 智慧ある者は足で歩き、愚かな者は手で歩く。

12. 「自分のことだけに気を取られなさい。この争いはあなたにとって関係ありません。」

13. 橋は砂糖で作られていない；一部の橋は風致がある。

14. 解けない謎は私には興味がない；これらのすべての謎は解けない。

15. ジョンは勤勉である；勤勉な者は誰もが幸せである。

16. 蛙は書籍を書かない；一部の人々は書籍を書く際にインクを使用する。

17. ポーカーは柔らかくない；すべての枕は柔らかい。

18. アンテロープは不格好ではない；優雅な動物は目を楽しませる。

19. 一部の叔父は大方でない；すべての商人は大方である。

20. 不幸な者は誰もが笑わない；幸せな者は誰もが嘆かない。

21. 聴覚に届く音楽は空気を振動させる；聴覚に届かない音楽は払う価値がない。

22. 彼は私に五ポンドを与えた；私は喜んだ。

23. 金髪のユダヤ人は太いミルラーではない；私のすべての友人は古いミルラーである。

24. 麦粉は食べ物として良い；オートミールは麦粉の一種である。

25. 一部の夢は恐ろしい；羊は誰もが恐ろしくない。

26. 富豪は通りで乞食にならない；すべての貧乏な者は帳簿をつけるべきである。

27. 盗賊は正直ではない；不正直な者の中には見つかった者がいる。

28. すべての蜂は不親切である；すべての仔犬は親切である。

29. 不確かな話は誰もが疑う；これらの話のどれも確かなものではない。

30. 「彼はあなたが去ったと言いました。」「彼は一つの真実な言葉も言わない。」

31. 彼の歌は一時間持続しない；一時間持続する歌は退屈である。

32. 嫁 Cake は誰もが健康でない；不健康な食事は避けるべきである。

33. 古い miser（小僧）は喜ばない；ある古い miser（小僧）は痩せている。

34. すべての duck（鴨）は歩く；何もが歩くものは優雅ではない。

35. すべての Professors（教授）は無知ではない；ある無知な人間は自慢高めている。

36. 牙痛は決して快適ではない；暖かさは決して不快ではない。

37. Bore（退屈者）は恐ろしい；あなたは bore（退屈者）である。

38. ある山は登れない；すべての stile（踏み切り）は登れる。

39. すべての Frenchmen（フランス人）は plumpudding（ポップコーン）を好まない；すべての Englishmen（イギリス人）は plumpudding（ポップコーン）を好む。

40. すべての idlers（怠け者）は名誉を得ない；ある画家は怠け者ではない。

41. すべての lobsters（カニ）は不合理ではない；合理的な生物は不可能を期待しない。

42. すべての kind deed（善行）は違法ではない；違法なことは恐れずにできる。

43. すべての fossils（化石）は愛に渡すことができない；どの oyster（牡蠣）も愛に渡すことができる。

44. 「これ以上耐えられない！」「では、これ以上耐えられないことが私にはかつて起こったことがない。」

45. すべての uneducated men（未教育者）は浅薄である；これらの学生はすべて教育を受けている。

46. すべての私の cousins（姉妹）は不正義である；すべての judges（裁判官）は不正義ではない。

47. 探査された国は龍にInfestedされていない；未探査された国は魅力的である。

48. すべての miser（小僧）は大方でない；ある古い人は大方でない。

49. 議事明哲な男は hyaenas（ハイエナ）を避ける；銀行員は議事不明哲ではない。

50. ある詩はoriginalである；すべてのoriginal作品は意に従って作成されるものではない。

51. すべての miser（小僧）は自愛的でない；misers（小僧）以外の人間は何も救えない。

52. すべての pale people（蒼白な人）はphlegmaticである；誰もがpaleでない人は詩的な外見を呈しない。

53. すべての spiders（蜘蛛）はwebを織る；あるcreature（生物）はwebを織らないにもかかわらず凶暴である。

54. 私のすべての cousins（姉妹）は正義でない；すべての judges（裁判官）は正義である。

55. Johnは勤勉である；勤勉な人間は決して不快ではない。

56. ウォーキング用のumbrella（傘）は有用である；旅程において無用なものには手を離すべきである。

57. あるpillow（枕）は柔らかい；何もがpoker（火鉢の棒）は柔らかくない。

58. 私は年老いて足が不自由である；年老いた商人は賭博家でない。

59. すべての eventful journey（興味深い旅）は決して忘れられない；uneventful journey（平凡な旅）は書籍に記す価値がない。

60. シュガーは甘い；あるsweet things（甘いもの）は子供たちによって好まれる。

61. Richardは怒り狂っている；誰もがRichard以外ではその馬に乗ることができない。

62. すべての jokes（冗談）は笑わすために作られている；何もが法案は冗談ではない。

63. 「私は新聞で見た。」「すべての新聞は嘘をついている。」

64. すべての nightmare（悪夢）は快適ではない；不快な経験は切望されない。

65. 議論を慎重に行う者は十分な零細金を持ち歩く；議論を粗忽に行う者は荷物を失う。

66. すべての蜂は不親切である；どの puppies も不親切ではない。

67. 彼は昨日ここに来た；彼は私の友人ではない。

68. すべての四足歩行動物は吹き出せない；いくつかの猫が四足歩行動物である。

69. 焙煎された肉は屠夫によって売られない；生肉は夕食には提供されない。

70. 金は重い；彼を静めることは金以外の何ものでもない。

71. いくつかの豚は野生である；どの豚も太っている。

72. すべての皇帝は歯科医ではない；すべての歯科医は子供たちに畏怖される。

73. すべての若くない者は歩きを好む；あなたも私は老いていない。

74. すべての刃物は鋭利である；いくつかの草が刃物である。

75. すべての独裁的な人物は人気を得ない；彼女は独裁的である。

76. いくつかの甘いものは不健康である；どの muffin も甘くない。

77. すべての軍人は詩を書かない；どの将校も平民ではない。

78. 誠に退屈な者は畏怖される；退屈な者は彼の滞在を延ばすために乞われることはない。

79. すべての猫頭鹰は満足させる；いくつかの理由が不満足である。

80. すべて私の堂兄弟は正義を持たない；すべての判事は公正である。

81. いくつかのパンは豊富である；すべてのパンは美味しい。

82. すべての薬草は味がしない；どのピルも薬効がある。

83. いくつかの授業は難解である；難解なものは注意を要する。

84. 予期された楽しみは私を不快にしない；あなたの訪問は予期されていない楽しみである。

85. カATERPILLARはeloquentでない；ジョーンズはeloquentである。

86. いくつかの薄毛の人々はウィッグを着用する；すべてのあなたの子供たちは髪を持っている。

87. すべての蜂は不親切である；不親切な生物は常に歓迎されない。

88. 破産者は富を持たない；いくつかの商人は破産者ではない。

89. ウサギは時々眠る；すべての動物は時々眠る。

90. 不適切に管理された企業は利益を生まない；鉄道は常に不適切に管理されない。

91. すべての人間が豚を見たことがある；誰も豚を賞賛しない。

\_\_\_\_\_\_\_\_\_\_\_\_\_\_

次のように前提を抽出し、結論（もし存在する場合）を導き出してください：

92. "ライオンは、私があなたが私が何回か追いかけられたことを知っているなら誰でも言えるように、非常に凶暴な動物である：そして、彼らのなかには特定の個体がいるかもしれないが、私はそれが普遍的な法則ではないと保証することはできないが、その個体たちはコーヒーを飲まない。"

前提1: ライオンは非常に凶暴である。
前提2: 特定のライオンはコーヒーを飲まない。

結論: あるライオンはコーヒーを飲まない。

93. "あなたがあなたが何か知恵があったら、古くから海に詳しい船員は粥を好まないことを知っていればよいのに！"

前提1: 古い船員は粥を好まない。
前提2: あなたは何か知恵がある。

結論: あなたの訪問は予期されていない楽しみである。

しかし、彼はあなたの伯父 Sir --"

「私の伯父だというのか！ 誤魔化するな！」

「伯父と呼べばいい。好きに呼びなさい。ただ、私には伯父が皆高齢者で、粥を好むだけのことは知っています！」

「では、あなたの伯父は--」

94. 「さあ、行こうじゃないか！これ以上この圧迫感に耐えられぬ。狭い店舗は快適とは言えないよ、それはよく分かっているだろう？」

「そうかい、出先で快適な気分を味わいたいか？」

「もちろん、私は快適な気分を味わいたい！そして、街の下の方には人混みの少ない店があると確信している。だから--」

95. 「医者たちは形而上学的なオルガン奏者ではないというので、あなたについて少しの事実が明らかになるよ、知っているでしょう？」

「どうしてそれを分かったのか？私をオルガンで演奏したことはないじゃないか？」

「いいえ、先生ですが、ブラウンの詩について話したことがあります。それから形而上学的な人物であることが分かりました。だから--」

96. 「話すな！私はあなたよりも多くの裕福な商人と知り合いだ：そして私には一つも古い老齢者ような貪欲な商人は存在しないと言える！」

「しかし、それはなぜ旧Mr. Brownに関係するのですか？」

「だって、彼は非常に裕福ではないか？」

「もちろん、そうよ。ではどうして？」

「だって、彼を貪欲な商人と呼ぶのは不自然だ！彼が商人でないかぎり、または彼が貪欲でない限り、そのような呼び方は適切ではない！」

97. 「あなたが尋ねてくれたことに感謝しているよ！私は今日とてもよく感じているんだ。」

「それは自然の恩恵か、技術の恩恵か？」

「技術だわ。医者は私に彼の特許薬を飲ませたんだ。」

「そうか、それなら医者の嘘つき呼ばわりはやめよう。少なくとも、その薬を飲んだ人は誰でも状態が良くなっている！」

98. 「いいえ、私はあなたを一目見て嫌いだ。そして、私の人形と遊ぶことにしよう。人形は決して不親切ではないわ。」

「つまり、あなたは人形よりも堂兄弟を好まないということか？あら、おばけのよう！」

「もちろんそうよ！堂兄弟は決して親切でない--少なくとも私が見た限りではね。」

「でも、それがあなたが証明しようとしていることなのか？もし堂兄弟が人形ではないと言っているなら、誰がそれを言ったのか？」

99. 「なぜガーネットについて話しているの？この距離から花を見分けることはできないよ！あなたはすべて赤い花だと言えるだけだ：望遠鏡も必要ないわ。」

「では、いくつかのベランダガーニュムは赤いでしょう？」

「認めます。でもそれから何が言いたいのですか？ あなたはきっと、そのうちのいくつかはベランダガーニュムだと言おうとしているでしょう！」

「もちろん、あなたが論理を追えるだけの知恵があるなら、そう言うべきです！しかし、何かを証明することに意味があるのか、私にはわかりませんよ？」

100. 「諸君、全体としてはまあ悪くない試験をパスしたね。今度は私が去る前に一言アドバイスをさせてもらうが、本当に学ぶことを望む者は努力するんだよ。」

「先生、学者たちの名代で感謝いたします！そして少なくとも一部だけでも真剣に学ぶことを望んでいる者がいることに誇りを感じます。」

「それはうれしいことだね。では、どうしてそれがそうなのか？」

「なぜなら、先生、私は彼らがどれほど努力しているか知っています──いくつかの子供たちについてです。誰がそれよりもよく知るでしょう？」

\_\_\_\_\_\_\_\_\_\_\_\_\_\_\_\_\_\_\_

次のスピーチから、サリュー（三段論法）や三段論法のような形式を持つ議論の一連を抽出し、その正確性を検証してください。

このスピーチは、友人が彼女が子供たちに対して授業にあまり熱心すぎるのかもしれませんと注意深く提案したとき、愛情深い母親が話すものとして想定されています。

101. 「彼らには自分たちの生き方がある。私たちはそれぞれに遺産を与えるわけにはいかない。そして、金は金額相当の価値がないと手に入らないことを汝も知っている通りである。したがって、彼らは働きなければならず、生活するためには働くしかない。しかし、何も知らないでどのように働けばよいのか？私に言わせれば、この時代では無知の場所はどこにもない！そして、すべての権威者は学ぶべき時期は若いうちであると一致して同意している。その後の記憶はほとんど価値がない。子供は一時間で大人が五時間で学べるよりも多くのことを学ぶことができる。したがって、学ぶ必要がある人々は、いつでも学ばない限り、若い頃に学ぶしかない。もちろん、その前に子供たちが健康であることが重要だという点については私は完全に同意する。では、医者は私に告げたのである。子供が健康を保つためには、一日の授業時間は六時間以内でなければならないし、週に少なくとも二日は半休を取るべきである。そして、その計画は私があなたに約束する通り、私たちが実行している計画である！私たちは六時間を超えることは決してないし、毎水曜日と土曜日に、彼らは午後一時以降授業をすることがない！だから、私の愛すべき子供たちの教育で私が冒すリスクをどのように想像できるのか、私は理解不能だと言わざるを得ない！」

終わり。


\backmatter

\chapter{訳者・数理校閲後記}

本書は、ルイス・キャロル（チャールズ・L・ドジソン）による知られざる論理学の傑作『論理の遊戯』（\textit{The Game of Logic}, 1887）の全4章を、Project Gutenberg（\#4763）の底本に基づき完全翻訳・数理校閲したものである。

キャロルが目指したのは、堅苦しい学問としての論理学ではなく、知的な歓びとユーモアに満ちた「遊戯（ゲーム）」としての記号論理であった。ダイアグラムの上で駒を動かしながら前提から結論を導き出す快感は、130年以上を経た現代においても全く色あせることはない。

本書が、論理的思考の純粋な楽しさと、アリスの作者が持っていた鋭利な数学的知性に触れる最良の契機となることを願う。

\begin{flushright}
2026年10月\\
教育情報分析研究会
\end{flushright}

\end{document}
